\documentclass[11pt,oneside]{article}

\usepackage[a4paper,margin=27mm,headheight=15pt]{geometry}
\usepackage[T1]{fontenc}
\usepackage[utf8]{inputenc}
\IfFileExists{libertinus.sty}{\usepackage{libertinus}}{\usepackage{lmodern}}
\usepackage{microtype}
\usepackage{amsmath,amssymb,amsthm,mathtools,mathrsfs}
\usepackage{aliascnt}
\usepackage{bm}
\usepackage{booktabs,longtable,array,tabularx}
\usepackage{graphicx}
\usepackage{pgfplots}
\usepackage{xcolor}
\usepackage{enumitem}
\usepackage{fancyhdr}
\usepackage{titlesec}
\usepackage{listings}
\usepackage{xurl}
\usepackage{hyperref}
\usepackage[nameinlink,noabbrev]{cleveref}

\definecolor{linkblue}{RGB}{25,75,135}
\definecolor{shadegray}{RGB}{246,247,249}
\definecolor{rulegray}{RGB}{195,201,208}
\pgfplotsset{compat=1.18}
\hypersetup{
  colorlinks=true,
  linkcolor=linkblue,
  citecolor=linkblue,
  urlcolor=linkblue,
  pdftitle={Fourier Decay of Rvachev's Up-Function: Final Synthesis},
  pdfsubject={Pointwise, metric, mean, extremal, and annular asymptotics of the dyadic sinc product},
  pdfkeywords={Rvachev up-function, Fabius function, sinc product, Fourier transform, Thue-Morse, transfer operator, asymptotic decay}
}

\pagestyle{fancy}
\fancyhf{}
\fancyhead[L]{\small Fourier decay of Rvachev's up-function}
\fancyhead[R]{}
\fancyfoot[C]{\thepage}
% Canonical project-wide mathematical notation.
% Canonical mathematical notation for active FabiusFunction LaTeX documents.
%
% This file is intentionally a plain TeX input rather than a LaTeX package.
% Every standalone document loads it by an explicit path relative to that
% document's directory; included fragments inherit the definitions from their
% root document.  Keep this file independent of document layout and metadata.
%
% Contract:
%   * one command has one mathematical meaning throughout the active corpus;
%   * names encode semantic roles rather than a document-local choice of letter;
%   * visually similar but differently normalized objects have different print;
%   * every command below is defined normatively in the notation catalogue.
%
% The active corpus excludes docs/archive/.  Historical sources may retain
% their original notation only inside an explicitly labelled quotation.

\makeatletter
\@ifundefined{mathbb}{%
  \PackageError{fabius-notation}{The amssymb package must be loaded first}%
    {Load amsmath and amssymb before inputting fabius-notation.tex.}%
}{}
\@ifundefined{operatorname}{%
  \PackageError{fabius-notation}{The amsmath package must be loaded first}%
    {Load amsmath before inputting fabius-notation.tex.}%
}{}
\makeatother

% Version stamp used by the catalogue and by static audits.
\newcommand{\FabiusNotationVersion}{2026-08-31}

% Number systems.  NaturalNumbers is explicitly the nonnegative convention.
\newcommand{\NaturalNumbers}{\mathbb{N}_{0}}
\newcommand{\PositiveIntegers}{\mathbb{N}_{>0}}
\newcommand{\IntegerNumbers}{\mathbb{Z}}
\newcommand{\RationalNumbers}{\mathbb{Q}}
\newcommand{\RealNumbers}{\mathbb{R}}
\newcommand{\ComplexNumbers}{\mathbb{C}}
\newcommand{\ExtendedRealNumbers}{\overline{\mathbb{R}}}

% The three principal Fabius--Rvachev functions.  Subscripts are deliberate:
% a bare F and a calligraphic F have several unrelated uses in the corpus.
\newcommand{\FabiusBounded}{F_{\mathrm{bd}}}
\newcommand{\FabiusGlobal}{F_{\mathrm{gl}}}
\newcommand{\FabiusQuantile}{Q_{F}}
\newcommand{\FabiusClampedQuantile}{Q_{F,\mathrm{cl}}}
\newcommand{\RvachevUp}{\operatorname{up}}

% Constants, elementary operators, and delimiters.
\newcommand{\EulerE}{\mathrm{e}}
\newcommand{\ImaginaryUnit}{\mathrm{i}}
\newcommand{\Differential}{\mathop{}\!\mathrm{d}}
\newcommand{\AbsoluteValue}[1]{\left\lvert #1\right\rvert}
\newcommand{\Norm}[1]{\left\lVert #1\right\rVert}
\newcommand{\Floor}[1]{\left\lfloor #1\right\rfloor}
\newcommand{\Ceiling}[1]{\left\lceil #1\right\rceil}
\newcommand{\FractionalPart}[1]{\left\{#1\right\}}
\newcommand{\PositivePart}[1]{\left(#1\right)_{+}}
\newcommand{\IndicatorOf}[1]{\mathbf{1}_{#1}}
\newcommand{\SetOf}[1]{\left\{#1\right\}}
\newcommand{\InnerProduct}[1]{\left\langle #1\right\rangle}
\newcommand{\CardinalityOf}[1]{\left\lvert #1\right\rvert}
\newcommand{\DefinitionEquals}{\mathrel{\coloneqq}}
\newcommand{\Sign}{\operatorname{sgn}}
\newcommand{\IdentityOperator}{\operatorname{id}}
\newcommand{\SpanOperator}{\operatorname{span}}
\newcommand{\TraceOperator}{\operatorname{tr}}
\newcommand{\RankOperator}{\operatorname{rank}}
\newcommand{\DiagonalOperator}{\operatorname{diag}}
\newcommand{\ResidueOperator}{\operatorname*{Res}}
\newcommand{\LeastCommonMultiple}{\operatorname{lcm}}
\newcommand{\OrderOperator}{\operatorname{ord}}
\newcommand{\PrincipalLogarithm}{\operatorname{Log}}
\newcommand{\FallingFactorial}[2]{\left(#1\right)^{\underline{#2}}}
\newcommand{\RisingFactorial}[2]{\left(#1\right)^{\overline{#2}}}

% Support, topology, convolution, and metric notation.
\newcommand{\SupportOperator}{\operatorname{supp}}
\newcommand{\SupportOf}[1]{\SupportOperator\!\left(#1\right)}
\newcommand{\TopologicalSupportOf}[1]{\operatorname{tsupp}\!\left(#1\right)}
\newcommand{\Convolution}{\mathbin{*}}
\newcommand{\MetricDistance}{\operatorname{dist}}
\newcommand{\TotalVariationDistance}{d_{\mathrm{TV}}}

% Probability.  Equality and convergence in law are not metric distance.
\newcommand{\Probability}{\mathbb{P}}
\newcommand{\Expectation}{\mathbb{E}}
\newcommand{\Variance}{\operatorname{Var}}
\newcommand{\Covariance}{\operatorname{Cov}}
\newcommand{\LawOperator}{\operatorname{Law}}
\newcommand{\LawOf}[1]{\LawOperator\!\left(#1\right)}
\newcommand{\EqualInLaw}{\mathrel{\overset{\mathrm{law}}{=}}}
\newcommand{\ConvergesInLaw}{\mathrel{\xrightarrow{\mathrm{law}}}}

% Fourier and integral transforms.  The printed labels expose normalization.
% FourierTwoPi uses exp(-2*pi*i*x*xi); FourierAngular uses exp(-i*x*omega).
\newcommand{\FourierTwoPi}[1]{\widehat{#1}_{\,2\pi}}
\newcommand{\FourierAngular}[1]{\widehat{#1}_{\,\mathrm{ang}}}
\newcommand{\LaplaceTransformSymbol}{\mathcal{L}}
\newcommand{\MellinTransformSymbol}{\mathcal{M}}
\newcommand{\LaplaceTransformOf}[1]{\LaplaceTransformSymbol\!\left[#1\right]}
\newcommand{\MellinTransformOf}[1]{\MellinTransformSymbol\!\left[#1\right]}
\newcommand{\MomentGeneratingFunction}[1]{M_{#1}}
\newcommand{\CharacteristicFunction}[1]{\varphi_{#1}}

% The two standard sinc functions must never share one symbol.
\newcommand{\SincRad}{\operatorname{sinc}_{\mathrm{rad}}}
\newcommand{\SincPi}{\operatorname{sinc}_{\pi}}

% Binary arithmetic and Thue--Morse notation.
\newcommand{\BinaryDigitSum}{\operatorname{wt}_{2}}
\newcommand{\ThueMorseSignSymbol}{\varepsilon}
\newcommand{\ThueMorseSign}[1]{\varepsilon_{#1}}
\newcommand{\PadicValuation}[1]{\nu_{#1}}
\newcommand{\TwoAdicValuation}{\nu_{2}}

% q-series notation.  Every base is an explicit argument.
\newcommand{\QPochhammer}[3]{\left(#1;#2\right)_{#3}}
\newcommand{\GaussianBinomial}[3]{%
  \genfrac{[}{]}{0pt}{}{#1}{#2}_{#3}%
}
\newcommand{\QInteger}[2]{\left[#1\right]_{#2}}

% Named special functions.
\newcommand{\Polylogarithm}{\operatorname{Li}}
\newcommand{\LambertW}{W}
\newcommand{\GammaFunction}{\Gamma}
\newcommand{\RiemannZeta}{\zeta}

% Asymptotics.  BigO and LittleO are operator tokens for existing formula
% grammar; prefer the At forms so that the limiting regime is printed.
\newcommand{\BigO}{\mathcal{O}}
\newcommand{\LittleO}{o}
\newcommand{\BigOAt}[2]{\mathcal{O}_{#2}\!\left(#1\right)}
\newcommand{\LittleOAt}[2]{o_{#2}\!\left(#1\right)}

\renewcommand{\headrulewidth}{0.4pt}

\titleformat{\section}{\Large\bfseries}{\thesection}{0.7em}{}
\titleformat{\subsection}{\large\bfseries}{\thesubsection}{0.7em}{}
\titleformat{\subsubsection}{\normalsize\bfseries}{\thesubsubsection}{0.7em}{}
\setcounter{secnumdepth}{3}
\setcounter{tocdepth}{2}
\setlist[itemize]{leftmargin=2em,itemsep=0.25em,topsep=0.4em}
\setlist[enumerate]{leftmargin=2.3em,itemsep=0.25em,topsep=0.4em}
\setlength{\emergencystretch}{2em}

\newtheorem{theorem}{Theorem}[section]

\newaliascnt{proposition}{theorem}
\newtheorem{proposition}[proposition]{Proposition}
\aliascntresetthe{proposition}

\newaliascnt{lemma}{theorem}
\newtheorem{lemma}[lemma]{Lemma}
\aliascntresetthe{lemma}

\newaliascnt{corollary}{theorem}
\newtheorem{corollary}[corollary]{Corollary}
\aliascntresetthe{corollary}

\newaliascnt{conjecture}{theorem}
\newtheorem{conjecture}[conjecture]{Conjecture}
\aliascntresetthe{conjecture}

\theoremstyle{definition}
\newaliascnt{definition}{theorem}
\newtheorem{definition}[definition]{Definition}
\aliascntresetthe{definition}

\newaliascnt{example}{theorem}
\newtheorem{example}[example]{Example}
\aliascntresetthe{example}

\theoremstyle{remark}
\newaliascnt{remark}{theorem}
\newtheorem{remark}[remark]{Remark}
\aliascntresetthe{remark}

\newaliascnt{warning}{theorem}
\newtheorem{warning}[warning]{Warning}
\aliascntresetthe{warning}

\crefname{theorem}{theorem}{theorems}
\Crefname{theorem}{Theorem}{Theorems}
\crefname{proposition}{proposition}{propositions}
\Crefname{proposition}{Proposition}{Propositions}
\crefname{lemma}{lemma}{lemmas}
\Crefname{lemma}{Lemma}{Lemmas}
\crefname{corollary}{corollary}{corollaries}
\Crefname{corollary}{Corollary}{Corollaries}
\crefname{conjecture}{conjecture}{conjectures}
\Crefname{conjecture}{Conjecture}{Conjectures}
\crefname{definition}{definition}{definitions}
\Crefname{definition}{Definition}{Definitions}
\crefname{example}{example}{examples}
\Crefname{example}{Example}{Examples}
\crefname{remark}{remark}{remarks}
\Crefname{remark}{Remark}{Remarks}
\crefname{warning}{warning}{warnings}
\Crefname{warning}{Warning}{Warnings}

\newcommand{\Phiup}{\Phi}
\newcommand{\spec}{\operatorname{spr}}
\newcommand{\esssup}{\operatorname*{ess\,sup}}

\newenvironment{proofidea}{\par\smallskip\noindent\textit{Proof architecture.}\ }{\hfill$\square$\par\smallskip}
\newenvironment{boxedremark}{%
  \par\smallskip\noindent\begin{center}\begin{minipage}{0.94\textwidth}%
  \color{black}\hrule\vspace{0.6em}\small
}{%
  \vspace{0.6em}\hrule\end{minipage}\end{center}\smallskip
}

\lstdefinestyle{code}{
  basicstyle=\ttfamily\footnotesize,
  backgroundcolor=\color{shadegray},
  frame=single,
  rulecolor=\color{rulegray},
  columns=fullflexible,
  keepspaces=true,
  showstringspaces=false,
  breaklines=true,
  xleftmargin=0.7em,
  xrightmargin=0.7em,
  aboveskip=0.8em,
  belowskip=0.8em
}

\title{\bfseries Fourier Decay of Rvachev's Up-Function\\[0.35em]
\large Final synthesis of pointwise, metric, mean, extremal, and annular asymptotics}
\author{Research article for Vladimir Reshetnikov's \texttt{ProveIt} project}
\date{26 August 2026}

\begin{document}
\maketitle

\begin{abstract}
Let
\[
 \Phi(x)=\widehat{\RvachevUp}(x)
 =\prod_{h=0}^{\infty}\SincRad\!\left(\frac{\pi x}{2^h}\right),
 \qquad \SincRad z=\frac{\sin z}{z},
\]
be the Fourier transform of Rvachev's compactly supported up-function.  This article audits
four independent first-wave analyses, the relevant lacunary-products paper of
Aistleitner--Hofer--Larcher, and the subsequent comparative audit, and then gives a single
self-contained synthesis.  The main conclusion is that the phrase ``the decay rate'' has no
unqualified meaning.  There is a universal deterministic factor
\[
 \exp\!\left(-\frac{(\log x)^2}{2\log2}\right),
\]
but the next power of $x$ depends on whether amplitude means a typical value, an arithmetic
mean, an RMS value, a pointwise upper envelope, or the largest value in a dyadic annulus.

An exact dyadic factorization reduces the remaining numerator to a lacunary sine cocycle for
the doubling map.  Its $L^q$ pressure yields a complete shell spectrum.  In particular, with
$E_\kappa(x)=x^{-\kappa}\exp(- (\log x)^2/(2\log2))$, the four principal powers are
\[
 \begin{aligned}
 \kappa_0&=3.151496129472319\ldots,&
 \kappa_1&=2.748070014871334\ldots,\\
 \kappa_2&=2.651496129472319\ldots,&
 \kappa_\infty&=2.359014879111741\ldots .
 \end{aligned}
\]
They correspond respectively to geometric mean / almost-everywhere dyadic rays, arithmetic
mean, root mean square, and the sharp pointwise upper envelope.  The old conjectural power
$\log_2(2\pi)$ is therefore exactly right for RMS decay and wrong for the other readings.
No universal power of $\log x$ repairs the mismatch.

The metric fluctuation is also settled.  For
$S_n(t)=\sum_{j<n}\log|2\sin(\pi2^jt)|$, an exact covariance computation gives
\[
 \int_0^1S_n(t)^2\,dt
 =\frac{\pi^2}{4}n-\frac{\pi^2}{3}(1-2^{-n}).
\]
Equivalently, $S_n$ is a coboundary perturbation of the reverse martingale-difference
sequence generated by $\log|\tan\pi t|$.  Hence the correct law-of-the-iterated-logarithm
constant is $\pi/2$ with denominator $\sqrt{2n\log\log n}$, or $\pi/\sqrt2$ with denominator
$\sqrt{n\log\log n}$.  The factor $\pi$ printed in the 2017 lacunary-products theorem is too
large by $\sqrt2$; the source of the discrepancy is the missing $1/2$ in the $L^2$ norm of
the cosine basis.

For the literal Cesaro question, the arithmetic-mean gauge
$g(x)=A_1E_{\kappa_1}(x)$, where $A_1=0.091266124131588\ldots$, gives shell means tending to
one and solves the requested two-sided boundedness condition.  It also gives convergence at
dyadic endpoints.  At unrestricted endpoints a nonconstant dyadic-mantissa profile remains,
governed by a continuous singular eigenmeasure of an explicit transfer operator.  We prove
that no gauge with a stabilizing within-shell shape (including powers and iterated logarithms,
regularly varying factors, and fixed log-periodic corrections) can remove this profile.  A completely
unrestricted exotic scale-dependent analytic gauge is not ruled out by that theorem.  The
multiplicatively natural logarithmic Cesaro mean does converge for all endpoints, with
normalizing constant $A_1^{\log}=0.09386063575678\ldots$.

Finally, the raw annular maximum
$\mathfrak M_k=\max_{2^k\le x\le2^{k+1}}|\Phi(x)|$ has a sharper boundary-layer law:
\[
 \log\mathfrak M_k
 =-\frac{\log2}{2}k^2-(\log2)\kappa_\infty k
 -\frac{W(c(k+1))^2+2W(c(k+1))}{2\log2}
 +O((\log\log k)^2),
\]
where $c=(\log2)\pi^{-1}\sqrt{3/2}$ and $W$ is the principal Lambert function.  Thus the
largest point in a shell approaches its left boundary and carries a second log-normal factor
on the $\log\log x$ scale.  In contrast, the unique lobe maxima immediately after powers of
two form explicit deep valleys with leading coefficient $-1/\log2$, twice the generic
log-square coefficient.  Together these results reconcile every substantive disagreement
among the source documents and give the final rigorous decay dictionary.
\end{abstract}

\tableofcontents
\newpage

\section{The question and the shape of the answer}
\label{sec:introduction}

Rvachev's up-function is the even $C^\infty$ probability density supported on $[-1,1]$
that is closely related to the Fabius distribution function.  Under the Fourier convention
\[
 \widehat f(\xi)=\int_{\RealNumbers}f(t)\EulerE^{-2\pi\ImaginaryUnit t\xi}\Differential t,
\]
its transform is the entire function
\begin{equation}
 \Phiup(z)=\widehat{\RvachevUp}(z)
 =\prod_{h=0}^{\infty}\SincRad\!\left(\frac{\pi z}{2^h}\right).
 \label{eq:Phi-definition}
\end{equation}
This formula is classical; see, for example, Arias de Reyna's systematic treatment
\cite{arias1982} and the Fabius/Rvachev development in \cite{proveit}.

The motivating questions in \cite{mse-extrema,mse-decay} ask for the amplitude of the
oscillations and, more specifically, for a positive elementary function $g$ such that an
average of $\AbsoluteValue{\Phiup}/g$ tends to one or at least stays between two positive constants.
A working draft in the \texttt{ProveIt} repository \cite{decay-draft} proposes a log-normal
factor together with the power $x^{-\log(2\pi)/\log2}$ and a possible power of $\log x$.
The log-normal factor is indeed fundamental, but the rest of the conclusion depends on
which norm or averaging operation is intended.

There are three separate obstructions to a single pointwise answer.

\begin{enumerate}
\item The transform has a zero at every nonzero integer, so no asymptotic
      $\AbsoluteValue{\Phiup(x)}\sim g(x)>0$ can hold on the full real ray.
\item After the deterministic denominator is removed, what remains is a Birkhoff product
      for the doubling map.  Its geometric mean, arithmetic mean, quadratic mean, and
      supremum have different exponential growth rates.
\item Even the unique maxima between consecutive zeros possess exceptional subsequences.
      The peaks associated with the binary period-two orbit attain the sharp global upper
      envelope, whereas the extrema in the first few intervals after a power of two are
      exponentially smaller on the $n^2$ scale.
\end{enumerate}

Thus the correct answer is a \emph{spectrum of decay rates}, not one formula.  A convenient
common scale is
\begin{equation}
 E_\kappa(x)
 :=\exp\!\left(-\frac{(\log x)^2}{2a}-\kappa\log x\right)
 =x^{-\kappa}\exp\!\left(-\frac{(\log x)^2}{2\log2}\right),
 \qquad a:=\log2.
 \label{eq:E-kappa}
\end{equation}
For fixed dyadic phase, the quadratic term in \eqref{eq:E-kappa} comes from the product of
the $n$ denominators $2^1,\dots,2^n$.  The linear term is where the lacunary numerator
enters.

\subsection{Main decay table}
\label{subsec:main-table}

The following table summarizes the results proved below.  In the finite-$q$ rows,
``normalized'' means the normalized $L^q$ norm on the dyadic shell
$[2^{n-1},2^n]$.

\begin{center}
\begin{tabularx}{\textwidth}{@{}>{\raggedright\arraybackslash}p{0.24\textwidth}
  >{\centering\arraybackslash}p{0.20\textwidth}
  >{\centering\arraybackslash}p{0.25\textwidth}
  X@{}}
\toprule
Amplitude notion & numerator scale $\Lambda$ & power $\kappa$ in $E_\kappa$ & conclusion\\
\midrule
Geometric mean ($q=0$) and typical dyadic ray
 & $1/2$
 & $\displaystyle \frac32+\frac{\log\pi}{a}$
 & shell geometric mean is exactly constant after normalization; the same exponent holds
   for Lebesgue-a.e. fixed phase\\[0.7em]
Arithmetic mean ($q=1$)
 & $\rho_1=0.6613226021\ldots$
 & $\displaystyle \frac{\log\pi+a/2-\log\rho_1}{a}$
 & shell means converge to a positive constant; ordinary Cesaro means have a nonconstant
   dyadic phase profile\\[0.7em]
Root mean square ($q=2$)
 & $2^{-1/2}$
 & $\displaystyle \frac{\log(2\pi)}{a}$
 & exact numerator identity $\int B_n^2=2^{-n}$; normalized shell RMS converges\\[0.7em]
Uniform envelope ($q=\infty$)
 & $\sqrt3/2$
 & $\displaystyle \frac32+\frac{\log(\pi/\sqrt3)}{a}$
 & global upper bound, sharp on $x=2^n(2/3)$; not a two-sided envelope for every extremum\\
\bottomrule
\end{tabularx}
\end{center}

The ordering is
\begin{equation}
 \begin{aligned}
 \kappa_\infty&=2.359014879\ldots<\kappa_2=2.651496129\ldots\\
 &<\kappa_1=2.748070014\ldots<\kappa_0=3.151496129\ldots .
 \end{aligned}
 \label{eq:kappa-order}
\end{equation}
Larger $q$ emphasizes rarer and larger values, hence gives a slower decay and a smaller
power $\kappa_q$.

\begin{boxedremark}
The coefficient $\log(2\pi)/\log2$ appearing in the working draft is not an accidental
near miss.  It is \emph{exactly} the root-mean-square power $\kappa_2$.  It is not the
arithmetic-mean power, the typical power, or the sharp uniform-envelope power.  No factor
$(\log x)^p$ can convert one of these powers into another, because their discrepancy is
already a nonzero multiple of $\log x$ in the logarithm of the amplitude.
\end{boxedremark}


\section{Audit map and final verdict}
\label{sec:audit-map}

The source corpus contains four first-wave documents and one second-wave comparative audit.
Their apparent disagreement is mostly a disagreement about which functional of $|\Phi|$ is
being estimated.  The following table records the final verdict before the detailed proofs.

\begin{center}
\small
\begin{tabularx}{\textwidth}{@{}>{\raggedright\arraybackslash}p{0.19\textwidth}X
 >{\raggedright\arraybackslash}p{0.31\textwidth}@{}}
\toprule
Source & Main contribution & Final assessment\\
\midrule
Wave 1, Document 1 \cite{wave1-doc1}
 & Exact dyadic reduction, typical law, martingale CLT/LIL, and a sharp pointwise envelope.
 & Substantively correct on these layers.  Its LIL constant $\pi/\sqrt2$ is the correct one;
   the argument is completed in \Cref{sec:fluctuations}.\\[0.45em]
Wave 1, Document 2 \cite{wave1-doc2}
 & Sharp raw dyadic-annulus maxima, Lambert-$W$ boundary layer, and explicit dyadic troughs.
 & Correct and strictly sharper than the other documents on annular maxima.  The proof is
   reconstructed in \Cref{sec:annular-maxima}.\\[0.45em]
Wave 1, Document 3 \cite{wave1-doc3}
 & Typical, extremal, arithmetic-mean, and $L^q$ interpretations; an analytic Cesaro gauge.
 & Correct except that it faithfully inherited the published LIL constant $\pi$, which is
   too large by $\sqrt2$.\\[0.45em]
Wave 1, Document 4 \cite{wave1-doc4}
 & Most systematic transfer-operator treatment, exact RMS identity, sharp envelope, phase
   profile, and fixed-offset valley theorem.
 & Correct modulo its explicitly imported Ruelle--Perron--Frobenius block; it omits the
   fluctuation and Lambert boundary-layer layers rather than contradicting them.\\[0.45em]
Aistleitner--Hofer--Larcher \cite{ahl2017}
 & Sharp metric study of the same lacunary products and a rigorous five-digit enclosure of
   the $L^1$ Perron root.
 & The order estimates and $L^1$ input are valid.  The printed LIL constant contains a
   factor-$\sqrt2$ normalization error, located precisely in \Cref{subsec:published-error}.\\[0.45em]
Comparative audit \cite{comparative-audit}
 & Reconciles the four readings, locates the LIL error, proves a stabilizing-gauge
   obstruction, and adds a logarithmic-mean repair.
 & Essentially correct.  Its phrase ``every elementary gauge'' must be read as the
   explicitly proved stabilizing/Hardy-field class; no theorem in the corpus excludes every
   conceivable scale-dependent analytic gauge.\\
\bottomrule
\end{tabularx}
\end{center}

The final answer can therefore be stated compactly.

\begin{theorem}[Decay dictionary]
\label{thm:decay-dictionary}
Let $E_\kappa$ be defined by \eqref{eq:E-kappa}.
\begin{enumerate}
\item For Lebesgue-a.e. dyadic phase, $|\Phi|=E_{\kappa_0}\exp(o(\log x))$,
with Gaussian fluctuations of variance $\pi^2/4$ per dyadic level and the corrected LIL of
\Cref{thm:correct-lil}.
\item The normalized shell $L^q$ means converge to positive constants for every fixed
$0<q<\infty$, with exponent $\kappa_q$ determined by the Perron root of
\eqref{eq:transfer-operator}.  In particular $q=1$ answers the literal mean-amplitude
question and $q=2$ gives $\kappa_2=\log_2(2\pi)$ exactly.
\item There is $C<\infty$ such that $|\Phi(x)|\le C E_{\kappa_\infty}(x)$ for $x\ge1$,
and equality up to one fixed constant holds on the ray $x=2^n(2/3)$.
\item The raw annular maxima have the additional Lambert-$W$ suppression stated in
\Cref{thm:annular-main}; the local lobe maxima have both sharp-envelope subsequences and the
fixed-offset dyadic valleys of \Cref{thm:dyadic-valleys}.
\item The gauge $A_1E_{\kappa_1}$ solves the two-sided Cesaro boundedness problem and the
dyadic-endpoint limit.  The unrestricted ordinary limit fails throughout the stabilizing
class of \Cref{thm:stabilizing-impossibility}; the logarithmic mean of
\Cref{thm:logarithmic-mean} converges without endpoint restrictions.
\end{enumerate}
\end{theorem}

\section{Exact products, zeros, and scaling laws}
\label{sec:exact-products}

\subsection{Entire convergence and three product forms}

The compact convergence of \eqref{eq:Phi-definition} follows from
\[
 \SincRad w=1-\frac{w^2}{6}+\BigO(w^4),
 \qquad
 \sum_{h\ge0}\frac{\AbsoluteValue z^2}{4^h}<\infty
\]
uniformly for $z$ in a compact set.  Hence $\Phiup$ is even, real on $\RealNumbers$, entire, and
$\Phiup(0)=1$.

Vieta's product $\SincRad t=\prod_{j\ge1}\cos(t/2^j)$ and Euler's sine product give two other
useful forms:
\begin{equation}
 \Phiup(z)
 =\prod_{j=1}^{\infty}\cos^j\!\left(\frac{\pi z}{2^j}\right)
 =\prod_{m=1}^{\infty}
   \left(1-\frac{z^2}{m^2}\right)^{1+\TwoAdicValuation(m)}.
 \label{eq:three-products}
\end{equation}
Indeed, in the cosine product the factor with denominator $2^j$ occurs for exactly $j$
pairs of indices.  In the Weierstrass product, the integer $m$ is a zero of the sinc factor
at level $h$ precisely when $2^h\mid m$, so its total multiplicity is $1+\TwoAdicValuation(m)$.

The elementary one-step scaling law is
\begin{equation}
 \Phiup(2x)=\SincRad(2\pi x)\Phiup(x),
 \qquad
 \Phiup(x)=\SincRad(\pi x)\Phiup(x/2).
 \label{eq:one-step-scaling}
\end{equation}

\subsection{Zeros and their multiplicities}

\begin{proposition}[Integer zero divisor]
\label{prop:zero-divisor}
The nonzero zeros of $\Phiup$ are exactly the nonzero integers.  The zero at
$m\in\IntegerNumbers\setminus\{0\}$ has multiplicity
\[
 d(m)=1+\TwoAdicValuation(\AbsoluteValue m).
\]
Moreover, if $b\ge1$ is an integer, $d=1+\TwoAdicValuation(b)$, and $z\uparrow b$, then
\begin{equation}
 \AbsoluteValue{\Phiup(z)}
 =b^{-d}\AbsoluteValue{\Phiup(b/2^d)}(b-z)^d\bigl(1+\BigO(b-z)\bigr).
 \label{eq:zero-local-coefficient}
\end{equation}
\end{proposition}

\begin{proof}
The divisor statement follows immediately from the last product in
\eqref{eq:three-products}.  For the local coefficient, the vanishing factors are those with
$0\le h\le\TwoAdicValuation(b)$.  At $z=b$ each has a simple zero and
\[
 \AbsoluteValue{\frac{\Differential}{\Differential z}\SincRad\!\left(\frac{\pi z}{2^h}\right)\Big|_{z=b}}
 =\frac1b.
\]
The product of the remaining factors is
\[
 \prod_{h=d}^{\infty}
 \AbsoluteValue{\SincRad\!\left(\frac{\pi b}{2^h}\right)}
 =\AbsoluteValue{\Phiup(b/2^d)}.
\]
Multiplication gives \eqref{eq:zero-local-coefficient}.
\end{proof}

\subsection{Exact dyadic renormalization}
\label{subsec:dyadic-renormalization}

Let
\[
 T(x)=2x\pmod1,
 \qquad
 s(x)=\AbsoluteValue{\sin(\pi x)},
 \qquad x\in[0,1],
\]
and define the lacunary product
\begin{equation}
 B_n(y):=\prod_{j=1}^{n}s(T^j y)
 =\prod_{j=1}^{n}\AbsoluteValue{\sin(\pi2^j y)}.
 \label{eq:Bny}
\end{equation}

\begin{theorem}[Exact dyadic factorization]
\label{thm:exact-dyadic-factorization}
For every $n\ge1$ and $y\in[1/2,1]$,
\begin{equation}
 \AbsoluteValue{\Phiup(2^n y)}
 =\AbsoluteValue{\Phiup(y)}\,\pi^{-n}y^{-n}2^{-n(n+1)/2}B_n(y).
 \label{eq:exact-factorization}
\end{equation}
For any real $\kappa$, set
\begin{equation}
 W_\kappa(y)
 :=\AbsoluteValue{\Phiup(y)}y^\kappa
   \exp\!\left(\frac{(\log y)^2}{2a}\right).
 \label{eq:W-kappa}
\end{equation}
Then the following normalized identity is exact:
\begin{equation}
 \frac{\AbsoluteValue{\Phiup(2^n y)}}{E_\kappa(2^n y)}
 =W_\kappa(y)
  \exp\!\left(n(\kappa a-\log\pi-a/2)\right)B_n(y).
 \label{eq:master-normalization}
\end{equation}
\end{theorem}

\begin{proof}
Iterating \eqref{eq:one-step-scaling} gives
\[
 \Phiup(2^n y)=\Phiup(y)\prod_{j=1}^{n}\SincRad(\pi2^j y).
\]
The product of the denominators is
$\pi^ny^n2^{1+\cdots+n}=\pi^ny^n2^{n(n+1)/2}$, proving
\eqref{eq:exact-factorization}.  Write $L=\log(2^ny)=na+\log y$ and divide
\eqref{eq:exact-factorization} by
$\exp(-L^2/(2a)-\kappa L)$.  The terms $-n\log y$ cancel, leaving
\eqref{eq:master-normalization}.
\end{proof}

Equation \eqref{eq:master-normalization} is the central identity of this article.  The
entire deterministic decay has been extracted into $E_\kappa$.  The bounded phase factor
$W_\kappa$ and the dynamical product $B_n$ contain all remaining information.

\subsection{The Thue--Morse polynomial}

Let $\tau_k=(-1)^{s_2(k)}$, where $s_2(k)$ is the number of ones in the binary expansion of
$k$.  The finite product identity
\begin{equation}
 P_n(x):=\prod_{j=0}^{n-1}(1-\EulerE^{2\pi\ImaginaryUnit2^jx})
 =\sum_{k=0}^{2^n-1}\tau_k\EulerE^{2\pi\ImaginaryUnit kx}
 \label{eq:TM-polynomial}
\end{equation}
follows by induction.  Taking moduli yields
\begin{equation}
 \AbsoluteValue{P_n(x)}=2^n\prod_{j=0}^{n-1}s(T^jx).
 \label{eq:TM-modulus}
\end{equation}
Thus the same products governing the decay of $\widehat{\RvachevUp}$ are the normalized moduli of
classical Thue--Morse trigonometric polynomials.  This is the bridge to Gelfond exponents,
Riesz products, and thermodynamic formalism; see \cite{gelfond,queffelec,fan-schmeling-shen}.

\subsection{A second exact scaling identity}
\label{subsec:boundary-identity}

The behavior just to the right of a large dyadic integer is exposed by a different exact
identity.

\begin{lemma}[Dyadic-boundary identity]
\label{lem:dyadic-boundary}
Let $N=2^{n-1}$ and $0<z<N$.  Then
\begin{equation}
 \AbsoluteValue{\Phiup(N+z)}
 =\frac{\AbsoluteValue{\Phiup(\tfrac12+z/(2N))}}
        {\AbsoluteValue{\Phiup(z/(2N))}}
  \left(\frac{z}{N+z}\right)^n\AbsoluteValue{\Phiup(z)}.
 \label{eq:dyadic-boundary}
\end{equation}
\end{lemma}

\begin{proof}
For $0\le h<n$, the number $N/2^h$ is an integer, so
\[
 \AbsoluteValue{\sin(\pi(N+z)/2^h)}=\AbsoluteValue{\sin(\pi z/2^h)}.
\]
The quotient of the first $n$ sinc factors is therefore $(z/(N+z))^n$ in absolute value.
Reindexing the remaining tails gives
\[
 \prod_{h=n}^{\infty}\SincRad\!\left(\frac{\pi(N+z)}{2^h}\right)
 =\Phiup\!\left(\frac12+\frac{z}{2N}\right),
 \quad
 \prod_{h=n}^{\infty}\SincRad\!\left(\frac{\pi z}{2^h}\right)
 =\Phiup\!\left(\frac{z}{2N}\right).
\]
Combining the finite and infinite parts proves the identity.
\end{proof}

\section{Exactly one extremum per unit interval}
\label{sec:extrema}

The Weierstrass product gives a short complete proof of the observed one-extremum property.

\begin{theorem}[Strict logarithmic concavity between zeros]
\label{thm:unique-extremum}
The function $\Phiup$ is positive and strictly decreasing on $(0,1)$.  For every integer
$m\ge1$, $\AbsoluteValue{\Phiup}$ has exactly one critical point
$\xi_m\in(m,m+1)$.  It is the strict maximum of $\AbsoluteValue{\Phiup}$ on that interval.  Moreover,
for every integer $m\ge0$,
\begin{equation}
 \Sign\Phiup(x)=(-1)^{s_2(m)},\qquad m<x<m+1.
 \label{eq:TM-sign}
\end{equation}
For $m\ge1$ the signed extremum is therefore
$\Phiup(\xi_m)=(-1)^{s_2(m)}M_m$, where
\[
 M_m:=\max_{m\le x\le m+1}\AbsoluteValue{\Phiup(x)}.
\]
\end{theorem}

\begin{proof}
Away from the integer zeros, termwise logarithmic differentiation of
\eqref{eq:three-products} is locally uniform and gives
\begin{align}
 \frac{\Differential}{\Differential x}\log\AbsoluteValue{\Phiup(x)}
 &=-2x\sum_{k=1}^{\infty}\frac{1+\TwoAdicValuation(k)}{k^2-x^2},
 \label{eq:log-derivative}\\
 \frac{\Differential^2}{\Differential x^2}\log\AbsoluteValue{\Phiup(x)}
 &=-2\sum_{k=1}^{\infty}(1+\TwoAdicValuation(k))
   \frac{k^2+x^2}{(k^2-x^2)^2}<0.
 \label{eq:strict-log-concavity}
\end{align}
For $0<x<1$, every denominator in \eqref{eq:log-derivative} is positive, so
$(\log\Phiup)'(x)<0$.  This proves strict decrease on the first interval.

Now let $m\ge1$.  The logarithmic derivative is strictly decreasing on $(m,m+1)$.
The pole contributed by $k=m$ forces it to $+\infty$ at the left endpoint, while the pole
contributed by $k=m+1$ forces it to $-\infty$ at the right endpoint.  It therefore has
exactly one zero, and strict log-concavity makes that point the unique maximum of the modulus.

Starting from $\Phiup(0)=1$, crossing the zero at $k$ changes the sign
$1+\TwoAdicValuation(k)$ times.  Therefore the sign on $(m,m+1)$ is
\[
 (-1)^{\sum_{k=1}^{m}(1+\TwoAdicValuation(k))}.
\]
Legendre's identity $\TwoAdicValuation(m!)=m-s_2(m)$ gives
\[
 \sum_{k=1}^{m}(1+\TwoAdicValuation(k))=m+\TwoAdicValuation(m!)=2m-s_2(m),
\]
whose parity is $s_2(m)$.  This proves \eqref{eq:TM-sign}.
\end{proof}

\begin{remark}
The critical points do \emph{not} in general converge to the midpoints of their unit
intervals.  In \Cref{sec:exceptional-extrema} we prove that, in every fixed interval offset
from a large power of two, the critical point approaches the \emph{right} endpoint at rate
$1/n$.
\end{remark}

\section{Dyadic-shell means and the transfer operator}
\label{sec:Lq-spectrum}

\subsection{Normalized shell means}

For $q\in(0,\infty)$ and a measurable $H$ on the shell
$D_n=[2^{n-1},2^n]$, define
\begin{equation}
 \mathcal M_{q,n}(H)
 :=\left(\frac1{2^{n-1}}\int_{2^{n-1}}^{2^n}\AbsoluteValue{H(x)}^q\Differential x\right)^{1/q}
 =\left(2\int_{1/2}^{1}\AbsoluteValue{H(2^ny)}^q\Differential y\right)^{1/q}.
 \label{eq:shell-Lq}
\end{equation}
The endpoint conventions are
\begin{align}
 \mathcal M_{0,n}(H)
 &:=\exp\!\left(\frac1{2^{n-1}}
       \int_{2^{n-1}}^{2^n}\log\AbsoluteValue{H(x)}\Differential x\right),
 \label{eq:shell-L0}\\
 \mathcal M_{\infty,n}(H)
 &:=\sup_{2^{n-1}\le x\le2^n}\AbsoluteValue{H(x)}.
 \label{eq:shell-Linf}
\end{align}
The logarithm in \eqref{eq:shell-L0} is locally integrable despite the finitely many zeros.

After $y=(x+1)/2$, the numerator in \eqref{eq:master-normalization} becomes
\begin{equation}
 \prod_{j=0}^{n-1}s(T^jx),
 \qquad 0\le x\le1.
 \label{eq:forward-product}
\end{equation}
This product is generated by an explicit Perron--Frobenius operator.

\subsection{The weighted doubling operator}

For $q>0$, define
\begin{equation}
 (\mathcal L_q f)(x)
 :=\frac12\left[
   \sin^q\!\left(\frac{\pi x}{2}\right)f(x/2)
  +\cos^q\!\left(\frac{\pi x}{2}\right)f((x+1)/2)
 \right].
 \label{eq:transfer-operator}
\end{equation}
The factor $1/2$ makes this the Lebesgue Perron operator with weight $s^q$.
For every integrable $f$,
\begin{equation}
 \int_0^1 f(x)\prod_{j=0}^{n-1}s(T^jx)^q\Differential x
 =\int_0^1(\mathcal L_q^n f)(x)\Differential x.
 \label{eq:transfer-duality}
\end{equation}

We use the following standard Ruelle--Perron--Frobenius fact.  The weights in
\eqref{eq:transfer-operator} are nonnegative rather than strictly positive at the two
endpoints, but the two-branch system is primitive, so the usual cone and
Lasota--Yorke proof applies without alteration; see \cite[Chs.~2--3]{baladi} and the
singular-potential discussion in \cite{fan1997,fan-schmeling-shen}.

\begin{proposition}[Perron root and spectral decomposition]
\label{prop:RPF}
For every $q>0$, $\mathcal L_q$ has a simple leading eigenvalue
$\rho_q>0$, a strictly positive eigenfunction $h_q$, and a positive eigenmeasure $\nu_q$,
which may be normalized by
\[
 \mathcal L_qh_q=\rho_qh_q,
 \qquad
 \mathcal L_q^*\nu_q=\rho_q\nu_q,
 \qquad
 \nu_q(h_q)=1.
\]
On a suitable Holder space $C^\alpha[0,1]$, with
$0<\alpha\le\min(1,q)$, there is a spectral gap and
\begin{equation}
 \rho_q^{-n}\mathcal L_q^nf
 =\nu_q(f)h_q+\BigO(\vartheta_q^n\Norm f_{C^\alpha}),
 \qquad 0<\vartheta_q<1.
 \label{eq:RPF-convergence}
\end{equation}
The same conclusion applies to nonnegative endpoint-power functions encountered below by
choosing a compatible Holder exponent.
\end{proposition}

\subsection{The full finite-\texorpdfstring{$q$}{q} decay spectrum}

Set
\begin{equation}
 \Lambda_q:=\rho_q^{1/q},
 \qquad
 \kappa_q:=\frac{\log\pi+a/2-\log\Lambda_q}{a},
 \qquad q>0.
 \label{eq:kappa-q}
\end{equation}

\begin{theorem}[Sharp finite-$q$ shell asymptotics]
\label{thm:Lq-spectrum}
For every $q\in(0,\infty)$ there is a constant $A_q\in(0,\infty)$ such that
\begin{equation}
 \mathcal M_{q,n}\!\left(\frac{\Phiup}{E_{\kappa_q}}\right)
 \longrightarrow A_q
 \qquad(n\to\infty).
 \label{eq:Lq-limit}
\end{equation}
More explicitly, with
\[
 \widetilde W_{q}(x)
 :=W_{\kappa_q}\!\left(\frac{x+1}{2}\right)^q,
\]
and the normalization in \Cref{prop:RPF},
\begin{equation}
 A_q^q=\left(\int_0^1h_q(x)\Differential x\right)\nu_q(\widetilde W_q).
 \label{eq:Aq}
\end{equation}
The number $\Lambda_q$ is nondecreasing in $q$, and therefore $\kappa_q$ is nonincreasing.
\end{theorem}

\begin{proof}
Raise \eqref{eq:master-normalization} to the $q$th power, use
\eqref{eq:shell-Lq}, and substitute $x=2y-1$.  This gives
\begin{align*}
 \mathcal M_{q,n}\!\left(\frac{\Phiup}{E_\kappa}\right)^q
 &=\exp\!\left(qn(\kappa a-\log\pi-a/2)\right)\\
 &\quad\times\int_0^1
 W_\kappa\!\left(\frac{x+1}{2}\right)^q
 \prod_{j=0}^{n-1}s(T^jx)^q\Differential x.
\end{align*}
By \eqref{eq:transfer-duality} and \eqref{eq:RPF-convergence}, the integral equals
\[
 \rho_q^n\left(\int_0^1h_q\right)
 \nu_q\!\left(W_\kappa((\,\cdot\,+1)/2)^q\right)(1+o(1)).
\]
The definition \eqref{eq:kappa-q} is exactly the cancellation condition
\[
 q(\kappa_q a-\log\pi-a/2)+\log\rho_q=0.
\]
This proves \eqref{eq:Lq-limit} and \eqref{eq:Aq}.

For each $n$, the $L^q$ norm of the product \eqref{eq:forward-product} on the probability
space $[0,1]$ is nondecreasing in $q$.  Taking $n$th roots and then the limit proves that
$\Lambda_q$ is nondecreasing.  Formula \eqref{eq:kappa-q} then gives the opposite
monotonicity of $\kappa_q$.
\end{proof}

\begin{remark}[No logarithmic-power correction]
\label{rem:no-log-correction}
For a fixed $q\in(0,\infty)$, replacing $E_{\kappa_q}$ by
$E_{\kappa_q}(\log x)^p$ multiplies the normalized shell mean by
$(na)^{-p}(1+o(1))$.  It therefore tends to $0$ when $p>0$ and to $+\infty$ when $p<0$.
At the sharp power, the required logarithmic exponent is $p=0$.
\end{remark}

\section{Geometric-mean and typical decay}
\label{sec:typical}

The $q=0$ endpoint is especially simple because Lebesgue measure is invariant under the
doubling map and
\begin{equation}
 \int_0^1\log\sin(\pi x)\Differential x=-\log2=-a.
 \label{eq:log-sine-integral}
\end{equation}

\begin{theorem}[Exact shell geometric mean]
\label{thm:geometric-mean}
Define
\begin{equation}
 \kappa_0:=\frac32+\frac{\log\pi}{a}
 =3.151496129472319\ldots .
 \label{eq:kappa0}
\end{equation}
Then
\begin{equation}
 \mathcal M_{0,n}\!\left(\frac{\Phiup}{E_{\kappa_0}}\right)
 =\exp\!\left(\int_0^1
    \log W_{\kappa_0}\!\left(\frac{x+1}{2}\right)\Differential x\right)
 \label{eq:exact-geometric-mean}
\end{equation}
for every $n\ge1$.  In particular, the normalized shell geometric mean is not merely
asymptotically constant; it is exactly independent of the shell.
\end{theorem}

\begin{proof}
Take logarithms in \eqref{eq:master-normalization}, substitute $x=2y-1$, and average.
The $n$ logarithmic sine terms each have integral $-a$ by invariance and
\eqref{eq:log-sine-integral}.  For \eqref{eq:kappa0},
\[
 \kappa_0a-\log\pi-a/2=a,
\]
so the linear contribution $na$ cancels the integral $-na$ exactly.
\end{proof}

The same exponent governs almost every fixed dyadic phase.

\begin{corollary}[Almost-everywhere dyadic-ray decay]
\label{cor:typical-ray}
For Lebesgue-a.e. $y\in[1/2,1]$,
\begin{equation}
 \log\AbsoluteValue{\Phiup(2^ny)}
 =-\frac{(\log(2^ny))^2}{2a}
  -\kappa_0\log(2^ny)+o(n).
 \label{eq:typical-ray}
\end{equation}
Equivalently,
\[
 \AbsoluteValue{\Phiup(2^ny)}=E_{\kappa_0}(2^ny)\EulerE^{o(n)}.
\]
\end{corollary}

\begin{proof}
The function $\log s$ is integrable and the doubling map is ergodic.  Birkhoff's theorem
therefore gives
\[
 \frac1n\log B_n(y)=\frac1n\sum_{j=1}^{n}\log s(T^jy)\longrightarrow-a
\]
for almost every $y$.  Substitute this in \eqref{eq:master-normalization} with
$\kappa=\kappa_0$.
\end{proof}

The $o(n)$ in \eqref{eq:typical-ray} is genuinely dynamical; it need not converge to a
constant.  Thus even the ``typical'' statement is a logarithmic asymptotic, not a pointwise
equivalent with one universal prefactor.


\section{Central limit theorem and the corrected iterated-logarithm law}
\label{sec:fluctuations}

The almost-everywhere exponent records only the mean of the logarithmic cocycle.  Its next
order is probabilistic, but in this particular problem the variance and the martingale
structure are completely explicit.  Put
\begin{equation}
 \psi(t):=\log|2\sin(\pi t)|,
 \qquad
 S_n(t):=\sum_{j=0}^{n-1}\psi(T^jt),
 \qquad Tt=2t\pmod1.
 \label{eq:psi-Sn}
\end{equation}
The exact centered logarithm in the alternative shell coordinates $x=2^k(1+t)$ is
\begin{equation}
 \log|\Phi(x)|
 =-\frac{(\log x)^2}{2a}-\kappa_0\log x+\Psi(1+t)+S_{k+1}(t),
 \label{eq:centered-log-final}
\end{equation}
where the mantissa term is explicitly
\begin{equation}
 \Psi(y):=\log W_{\kappa_0}(y/2),\qquad 1<y<2.
 \label{eq:Psi-mantissa}
\end{equation}
It is bounded on every compact subinterval of $(1,2)$.  Thus all metric fluctuations are
those of $S_n$.

\subsection{Exact covariance and finite-\texorpdfstring{$n$}{n} variance}

The classical Fourier expansion
\begin{equation}
 \psi(t)=-\sum_{m=1}^{\infty}\frac{\cos(2\pi mt)}{m}
 \quad\text{in }L^2(0,1)
 \label{eq:clausen}
\end{equation}
contains the whole calculation.

\begin{theorem}[Exact variance]
\label{thm:exact-variance}
The function $\psi$ has mean zero and, for every $n\ge1$,
\begin{equation}
 \int_0^1 S_n(t)^2\Differential t
 =\frac{\pi^2}{4}n-\frac{\pi^2}{3}\bigl(1-2^{-n}\bigr).
 \label{eq:variance-exact}
\end{equation}
Consequently the asymptotic variance per dyadic level is
\begin{equation}
 \sigma^2=\lim_{n\to\infty}\frac1n\int_0^1S_n^2=\frac{\pi^2}{4}.
 \label{eq:sigma}
\end{equation}
\end{theorem}

\begin{proof}
The mean is zero by the classical log-sine integral.  Parseval and
$\int_0^1\cos^2(2\pi mt)\Differential t=1/2$ give
\[
 C_0:=\int_0^1\psi(t)^2\Differential t
 =\frac12\sum_{m\ge1}\frac1{m^2}=\frac{\pi^2}{12}.
\]
For $r\ge1$, orthogonality leaves only frequencies $m=2^r\ell$:
\[
 C_r:=\int_0^1\psi(t)\psi(T^rt)\Differential t
 =\frac12\sum_{\ell\ge1}\frac1{2^r\ell}\frac1\ell
 =\frac{\pi^2}{12\,2^r}.
\]
Stationarity now yields
\[
 \int S_n^2=nC_0+2\sum_{r=1}^{n-1}(n-r)C_r.
\]
Since $\sum_{r=1}^{n-1}(n-r)2^{-r}=n-2+2^{1-n}$, this is exactly
\eqref{eq:variance-exact}.
\end{proof}

\subsection{A one-line Gordin decomposition}

Let $\mathcal P$ be the unweighted Perron operator of the doubling map,
\[
 (\mathcal P f)(x)=\frac12\left(f(x/2)+f((x+1)/2)\right).
\]
From \eqref{eq:clausen}, $\mathcal P\psi=\psi/2$.  Therefore
\begin{equation}
 d:=2\psi-\psi\circ T=\log|\tan(\pi t)|,
 \qquad \mathcal P d=0,
 \label{eq:martingale-difference}
\end{equation}
and telescoping gives the exact identity
\begin{equation}
 S_n=\sum_{j=0}^{n-1}d\circ T^j-\psi+\psi\circ T^n.
 \label{eq:Gordin-exact}
\end{equation}
The sequence $(d\circ T^j)$ is a stationary ergodic reverse martingale-difference sequence
for the decreasing filtration $(T^{-j}\mathcal B)$.  Moreover
\begin{equation}
 \int_0^1d(t)^2\Differential t
 =\frac2\pi\int_0^{\pi/2}\log^2(\tan u)\Differential u
 =\frac{\pi^2}{4}.
 \label{eq:d-variance}
\end{equation}
The final integral follows from the standard log-sine and log-cosine integrals; it also
agrees with \eqref{eq:sigma}, as it must.

\begin{theorem}[CLT and corrected LIL]
\label{thm:correct-lil}
With respect to Lebesgue measure on $t\in(0,1)$,
\begin{equation}
 \frac{S_n}{\sqrt n}\ \Longrightarrow\
 N\!\left(0,\frac{\pi^2}{4}\right).
 \label{eq:CLT}
\end{equation}
For Lebesgue-a.e. $t$,
\begin{equation}
 \limsup_{n\to\infty}\frac{S_n(t)}{\sqrt{2n\log\log n}}=\frac\pi2,
 \qquad
 \liminf_{n\to\infty}\frac{S_n(t)}{\sqrt{2n\log\log n}}=-\frac\pi2.
 \label{eq:LIL-correct}
\end{equation}
Equivalently, with denominator $\sqrt{n\log\log n}$, the two constants are
$\pm\pi/\sqrt2$.
\end{theorem}

\begin{proof}
The tail estimate $\operatorname{Leb}\{|d|>u\}=O(\EulerE^{-u})$ shows that
$d\in L^p(0,1)$ for every finite $p$.  With
$\DyadicSigmaField_j=T^{-j}\mathcal B$, the variables $d\circ T^j$ form a stationary reverse
martingale-difference sequence.  Moreover,
\[
 \frac1n\sum_{j=0}^{n-1}
 \mathbb E\!\left[(d\circ T^j)^2\mid\DyadicSigmaField_{j+1}\right]
 =\frac1n\sum_{j=0}^{n-1}(\mathcal P d^2)\circ T^{j+1}
 \longrightarrow\int_0^1d^2
\]
almost surely by the ergodic theorem, and the conditional Lindeberg condition follows from
$d\in L^2$.  The stationary martingale CLT and the martingale LIL therefore apply to
$\sum_{j<n}d\circ T^j$; see Stout \cite{stout1970} and Heyde--Scott
\cite{heyde-scott1973}.  The variance is \eqref{eq:d-variance}.  It remains only to remove
the two coboundary terms in \eqref{eq:Gordin-exact}.  The positive part of $\psi$ is bounded
by $\log2$, while near the integers
\[
 \operatorname{Leb}\{\psi<-u\}=O(\EulerE^{-u}).
\]
Hence
$\sum_n\operatorname{Leb}\{|\psi\circ T^n|>\varepsilon\sqrt{n\log\log n}\}<\infty$.
Borel--Cantelli makes the boundary terms $o(\sqrt{n\log\log n})$ almost surely.  Replacing
the threshold by $\varepsilon\sqrt n$ also makes them $o(\sqrt n)$ in probability, which
completes the CLT.
\end{proof}

In terms of the Fourier transform, for a.e. fixed $y=1+t\in(1,2)$ and
$x_k=2^ky$, \eqref{eq:centered-log-final} and \eqref{eq:LIL-correct} give
\begin{equation}
 \limsup_{k\to\infty}
 \frac{\log|\Phi(x_k)|+(\log x_k)^2/(2a)+\kappa_0\log x_k-\Psi(y)}
 {\sqrt{2(k+1)\log\log(k+1)}}=\frac\pi2,
 \label{eq:LIL-Phi}
\end{equation}
with the corresponding negative liminf.  The multiplicative fluctuation is therefore
\[
 \exp\!\left(\pm\left(\frac\pi{\sqrt2}+o(1)\right)
 \sqrt{\frac{\log x}{\log2}\,\log\log\log x}\right),
\]
which is subpolynomial in $x$ but larger than every fixed power of $\log x$.

\subsection{The normalization error in the published theorem}
\label{subsec:published-error}

Theorem 1.2 of Aistleitner--Hofer--Larcher \cite{ahl2017} prints the upper and recurrent
lower constants as $\pi$ when the denominator is $\sqrt{n\log\log n}$.  That value is
incompatible with the exact variance \eqref{eq:variance-exact}.  The paper's own variance
calculation effectively uses
$\int_0^1\cos^2(2\pi mt)\Differential t=1$ rather than $1/2$, doubling $\sigma^2$ from
$\pi^2/4$ to $\pi^2/2$.  Restoring the missing factor gives
$\sigma=\pi/2$ and therefore the constant $\sqrt2\sigma=\pi/\sqrt2$ in the
$\sqrt{n\log\log n}$ normalization.  This correction changes no exponent or qualitative
conclusion in that paper, but it matters for the sharp constant.

\section{The exact RMS power}
\label{sec:RMS}

For $q=2$, the transfer operator simplifies dramatically:
\begin{equation}
 \mathcal L_2\mathbf1
 =\frac12\left(\sin^2\frac{\pi x}{2}+\cos^2\frac{\pi x}{2}\right)
 =\frac12\mathbf1.
 \label{eq:L2-one}
\end{equation}
Its operator norm on $C[0,1]$ is at most $1/2$, so
\begin{equation}
 \rho_2=\frac12,
 \qquad
 \Lambda_2=\frac1{\sqrt2}.
 \label{eq:rho2}
\end{equation}
There is also a completely finite Fourier proof.  Parseval applied to
\eqref{eq:TM-polynomial} gives
\[
 \int_0^1\AbsoluteValue{P_n(x)}^2\Differential x=2^n,
\]
while \eqref{eq:TM-modulus} gives
\begin{equation}
 \int_0^1\prod_{j=0}^{n-1}s(T^jx)^2\Differential x=2^{-n}.
 \label{eq:exact-L2-product}
\end{equation}

\begin{corollary}[RMS decay]
\label{cor:RMS-decay}
The sharp root-mean-square exponent is
\begin{equation}
 \boxed{\displaystyle
 \kappa_2=\frac{\log(2\pi)}{\log2}
 =2.651496129472318\ldots .}
 \label{eq:kappa2}
\end{equation}
There is a constant $A_2\in(0,\infty)$ such that
\begin{equation}
 \mathcal M_{2,n}\!\left(\frac{\Phiup}{E_{\kappa_2}}\right)
 \longrightarrow A_2.
 \label{eq:RMS-limit}
\end{equation}
No power of $\log x$ belongs in the leading RMS normalization.
\end{corollary}

\begin{remark}[What the draft computed]
The power $x^{-\log(2\pi)/\log2}$ in \cite{decay-draft} is precisely
$x^{-\kappa_2}$.  Thus the draft identified the quadratic-mean scale, albeit without
specifying that norm.  Its claimed pointwise equivalent and its proposed logarithmic
correction do not follow and are incompatible with the exact shell identity
\eqref{eq:exact-L2-product}, the integer zeros, and the exceptional extrema proved later.
\end{remark}

\section{Arithmetic-mean decay and the Cesaro question}
\label{sec:L1-Cesaro}

\subsection{The Perron root \texorpdfstring{$\rho_1$}{rho1}}

For $q=1$,
\begin{equation}
 (\mathcal L_1f)(x)
 =\frac12\left[
  \sin\frac{\pi x}{2}\,f(x/2)
 +\cos\frac{\pi x}{2}\,f((x+1)/2)
 \right].
 \label{eq:L1}
\end{equation}
There appears to be no elementary closed form for its Perron root.  It is nevertheless an
exact, canonical constant, and it is easy to compute stably:
\begin{equation}
 \rho_1=0.661322602060565\ldots .
 \label{eq:rho1-numerical}
\end{equation}
The corresponding power is
\begin{equation}
 \boxed{\displaystyle
 \kappa_1
 =\frac{\log\pi+a/2-\log\rho_1}{a}
 =2.748070014871334\ldots .}
 \label{eq:kappa1}
\end{equation}

The decimal in \eqref{eq:rho1-numerical} is numerical, but the following elementary
Collatz--Wielandt certificate provides a rigorous enclosure.

\begin{proposition}[A rational enclosure for $\rho_1$]
\label{prop:rho1-rigorous-bound}
One has
\begin{equation}
 0.66126807<\rho_1<0.66134891,
 \label{eq:rho1-rigorous-bound}
\end{equation}
and consequently
\begin{equation}
 2.74801262<\kappa_1<2.74818899.
 \label{eq:kappa1-rigorous-bound}
\end{equation}
\end{proposition}

\begin{proof}
For a positive continuous test function $f$, the Collatz--Wielandt inequalities give
\[
 \inf_{x\in[0,1]}\frac{\mathcal L_1f(x)}{f(x)}
 \le\rho_1\le
 \sup_{x\in[0,1]}\frac{\mathcal L_1f(x)}{f(x)}.
\]
Take
\[
 f(x)=F(\sin\pi x),\qquad
 F(u)=1+\frac{376189}{10^6}u-\frac{69093}{10^6}u^2
       +\frac{15483}{10^6}u^3.
\]
This function is positive on $[0,1]$.  Put $t=\tan(\pi x/4)\in[0,1]$,
$u=2t/(1+t^2)$, and $v=(1-t^2)/(1+t^2)$.  Then
\[
 \frac{\mathcal L_1f(x)}{f(x)}
 =\frac{uF(u)+vF(v)}{2F(2uv)},
\]
a rational function of $t$.  Exact Sturm-sequence checks show that the numerators of its
differences from the two rational endpoints in \eqref{eq:rho1-rigorous-bound} are positive
and have no zeros on $[0,1]$.  A compact exact-arithmetic verification is included in
\Cref{app:certificate}.
\end{proof}

\begin{remark}[Published enclosure]
Aistleitner--Hofer--Larcher prove the sharper rigorous interval
\begin{equation}
 0.66130<\rho_1<0.66135,
 \label{eq:rho1-AHL}
\end{equation}
by the Fouvry--Mauduit transfer recursion.  It implies
$2.74801024<\kappa_1<2.74811933$.  The wider interval above is retained because its
certificate is short, exact, and executable from the appendix alone.
\end{remark}

\subsection{Sharp dyadic-shell mean}

By \Cref{thm:Lq-spectrum}, there is $A_1>0$ such that
\begin{equation}
 2\int_{1/2}^{1}
 \frac{\AbsoluteValue{\Phiup(2^ny)}}{E_{\kappa_1}(2^ny)}\Differential y
 \longrightarrow A_1.
 \label{eq:L1-shell-mean}
\end{equation}
Therefore the choice
\begin{equation}
 g_1(x):=A_1E_{\kappa_1}(x)
 \label{eq:g1}
\end{equation}
makes the mean of $\AbsoluteValue{\Phiup}/g_1$ on each dyadic shell tend to one.  Chebyshev
collocation of the transfer operator, followed by the spectral formula
\eqref{eq:Aq}, gives
\begin{equation}
 A_1=0.091266124131588\ldots .
 \label{eq:A1-numeric}
\end{equation}
The decimal is numerical; the exact definition is the positive spectral expression
\eqref{eq:Aq}.  This gives the precise answer to the natural shell version of the question
in \cite{mse-decay}.

\begin{proposition}[Regularity of the arithmetic-mean gauge]
\label{prop:gauge-regularity}
The function $g_1(x)=A_1E_{\kappa_1}(x)$ is positive and real-analytic on $(0,\infty)$ and
strictly decreasing for $x\ge1$.  For every fixed integer $m\ge0$ there is $X_m$ such that
\begin{equation}
 (-1)^m g_1^{(m)}(x)>0,\qquad x\ge X_m.
 \label{eq:eventual-alternating-derivatives}
\end{equation}
Thus it has every one of the eventual regularity properties requested in the original
question, although the threshold may depend on the derivative order.
\end{proposition}

\begin{proof}
Put $u=\log x$ and $G(u)=g_1(\EulerE^u)$.  Then
\[
 G(u)=A_1\exp\!\left(-\frac{u^2}{2a}-\kappa_1u\right),
 \qquad
 x^m\frac{\Differential^m}{\Differential x^m}
 =D_u(D_u-1)\cdots(D_u-m+1).
\]
Writing $p(u)=u/a+\kappa_1$, repeated differentiation gives
\[
 (-1)^m\frac{x^mg_1^{(m)}(x)}{g_1(x)}
 =p(u)^m+O_m(p(u)^{m-1})
 \qquad(u\to\infty).
\]
The right-hand side is positive for all sufficiently large $u$.  The case $m=1$ also follows
directly from $g_1'(x)/g_1(x)=-(\log x/a+\kappa_1)/x<0$ for $x\ge1$.
\end{proof}

It also answers the weaker full-ray boundedness question.

\begin{corollary}[Two-sided Cesaro boundedness]
\label{cor:Cesaro-boundedness}
For every fixed $t_0>0$, there are constants $0<c<C<\infty$ such that, for all sufficiently
large $t$,
\begin{equation}
 c<\frac1{t-t_0}\int_{t_0}^{t}
 \frac{\AbsoluteValue{\Phiup(x)}}{E_{\kappa_1}(x)}\Differential x<C.
 \label{eq:Cesaro-boundedness}
\end{equation}
After multiplying $E_{\kappa_1}$ by $A_1$, the same average tends to one along the dyadic
endpoints $t=2^n$.
\end{corollary}

\begin{proof}
Equation \eqref{eq:L1-shell-mean} says that the integral over the shell
$[2^{n-1},2^n]$ is asymptotic to $A_1 2^{n-1}$.  Summing the geometric sequence of completed
shells proves the dyadic-endpoint limit.  If $t$ lies inside the next shell, the completed
shells already contribute a positive constant times $t$, while the entire current shell
contributes at most another constant times $t$.  This gives \eqref{eq:Cesaro-boundedness}.
\end{proof}

\subsection{Why the ordinary Cesaro limit has a dyadic phase}
\label{subsec:phase-profile}

The full limit in the stronger condition of \cite{mse-decay} does not follow from the shell
limit.  In fact, the natural normalization \eqref{eq:g1} has a nonconstant phase profile.
This can be stated exactly in terms of the eigenmeasure of $\mathcal L_1$.

Normalize $h_1,\nu_1$ as in \Cref{prop:RPF}, and put
\[
 f_1(x):=W_{\kappa_1}\!\left(\frac{x+1}{2}\right),
 \qquad
 \mathscr F(z):=
 \frac{\nu_1(f_1\mathbf1_{[0,z]})}{\nu_1(f_1)},
 \quad 0\le z\le1.
 \label{eq:phase-distribution}
\]
The measure $f_1\,\Differential\nu_1$ is non-atomic, so $\mathscr F$ is continuous.  Here is a
short verification.  The endpoint atoms of $\nu_1$ vanish because the corresponding
transfer weights vanish.  For $z\in(0,1)\setminus\{1/2\}$, the eigenmeasure equation gives
\[
 \rho_1\nu_1(\{z\})=\frac12s(z)\nu_1(\{Tz\}).
\]
(The possible atom at $1/2$ is fed only by the already vanishing endpoint atoms.)  Hence any
remaining atom would generate forward atoms whose masses grow by at least the factor
$2\rho_1>1$ at every step.  A periodic orbit gives an immediate contradiction, and a
nonperiodic orbit would produce masses exceeding the finite total mass.  Thus $\nu_1$, and
therefore $f_1\,\Differential\nu_1$, has no atoms.

\begin{theorem}[Dyadic phase limit]
\label{thm:phase-limit}
For every $y\in[1/2,1]$ and every fixed $t_0>0$,
\begin{equation}
 \lim_{n\to\infty}
 \frac1{2^ny-t_0}\int_{t_0}^{2^ny}
 \frac{\AbsoluteValue{\Phiup(x)}}{A_1E_{\kappa_1}(x)}\Differential x
 =\mathscr P(y):=\frac{1+\mathscr F(2y-1)}{2y}.
 \label{eq:phase-limit}
\end{equation}
The continuous function $\mathscr P$ satisfies
$\mathscr P(1/2)=\mathscr P(1)=1$ but is not identically one.  Consequently the ordinary
Cesaro limit for the natural elementary scale $A_1E_{\kappa_1}$ does not exist.
\end{theorem}

\begin{proof}
The spectral decomposition gives, for $z\in[0,1]$,
\[
 \rho_1^{-n}\int_0^z f_1(x)
 \prod_{j=0}^{n-1}s(T^jx)\Differential x
 \longrightarrow
 \left(\int_0^1h_1\right)\nu_1(f_1\mathbf1_{[0,z]}).
\]
Approximation from above and below by Holder functions justifies the indicator; the
leading conformal measure has no atoms.  Returning to the shell variable, the partial
integral from $2^{n-1}$ to $2^ny$ is asymptotic to
\[
 2^{n-1}A_1\mathscr F(2y-1).
\]
The sum of all preceding complete shells is asymptotic to $2^{n-1}A_1$.  Division by
$2^nyA_1$ gives \eqref{eq:phase-limit}.

It remains to prove that $\mathscr F$ is not the identity distribution.  We first note that
$\nu_1$ is singular with respect to Lebesgue measure.  Otherwise its nonzero absolutely
continuous part would yield an $L^1$ density $r\ge0$ satisfying
\begin{equation}
 \rho_1r(x)=s(x)r(Tx)
 \quad\text{for a.e. }x.
 \label{eq:conformal-density}
\end{equation}
On the positive invariant support of $r$, iteration gives
\[
 r(T^nx)=\frac{\rho_1^nr(x)}{\prod_{j=0}^{n-1}s(T^jx)}.
\]
For Lebesgue-a.e. $x$, Birkhoff's theorem makes the logarithmic growth rate of the right
side $\log(2\rho_1)>0$.  On the other hand, for any $r\in L^1$ and any $\epsilon>0$,
Markov's inequality and Borel--Cantelli imply
$r(T^nx)\le\EulerE^{\epsilon n}$ eventually for almost every $x$.  Choosing
$0<\epsilon<\log(2\rho_1)$ is a contradiction.  Thus $\nu_1$ is singular, and multiplying
it by the continuous function $f_1$ leaves it singular.  Hence its normalized distribution
$\mathscr F$ cannot equal $z$, the distribution function of Lebesgue measure.  But
$\mathscr P\equiv1$ would force
$1+\mathscr F(2y-1)=2y$, or $\mathscr F(z)=z$.  Therefore $\mathscr P$ is nonconstant.
\end{proof}

For orientation, the following finite-level computation displays the nonconstant profile in
terms of the shell mantissa $m=2y\in[1,2]$.  It uses the normalized density
$f_1(z)\prod_{j<18}s(T^jz)$ on a dyadic midpoint grid; the theorem, not the plot, supplies
the limiting statement.

\begin{center}
\begin{tikzpicture}
\begin{axis}[
 width=0.82\textwidth,
 height=0.34\textwidth,
 xlabel={dyadic mantissa $m$},
 ylabel={$\mathscr P(m/2)$},
 xmin=1,xmax=2,
 ymin=0.90,ymax=1.13,
 grid=major,
 tick align=outside,
]
\addplot+[mark=*] coordinates {
 (1.0000,1.000000000)
 (1.0625,0.945118496)
 (1.1250,0.925107303)
 (1.1875,0.961711083)
 (1.2500,0.971580861)
 (1.3125,0.995925058)
 (1.3750,1.079490550)
 (1.4375,1.114228254)
 (1.5000,1.086316716)
 (1.5625,1.058951765)
 (1.6250,1.076832568)
 (1.6875,1.108201849)
 (1.7500,1.099910476)
 (1.8125,1.079553278)
 (1.8750,1.062385735)
 (1.9375,1.032028829)
 (2.0000,1.000000000)
};
\end{axis}
\end{tikzpicture}
\end{center}

\begin{remark}
A slowly varying correction, including every power of $\log x$, cannot remove this phase.
For a slowly varying $L$, the ratio $L(2^ny)/L(2^n)$ tends uniformly to one for
$y\in[1/2,1]$, so the profile \eqref{eq:phase-limit} survives.  An exact all-phase
normalizer would have to encode the singular conformal distribution, not merely add a
smooth logarithmic factor.
\end{remark}


\subsection{How far the unrestricted Cesaro obstruction extends}
\label{subsec:stabilizing-gauges}

The nonconstant profile in \Cref{thm:phase-limit} is not an artifact of having omitted a
power of $\log x$.  It survives every gauge whose relative shape settles down inside
successive dyadic shells.

\begin{theorem}[Impossibility in the stabilizing class]
\label{thm:stabilizing-impossibility}
Let $g:(X_0,\infty)\to(0,\infty)$ be a gauge.  Suppose there are constants $c_k>0$ and a
continuous $u:[0,1]\to(0,\infty)$ such that
\begin{equation}
 \frac{g(2^k(1+z))}{c_kE_{\kappa_1}(2^k(1+z))}
 \longrightarrow u(z)
 \quad\text{uniformly for }0\le z\le1.
 \label{eq:stabilizing-shape}
\end{equation}
Then the unrestricted limit
\begin{equation}
 \frac1{t-X_0}\int_{X_0}^t\frac{|\Phi(x)|}{g(x)}\Differential x\longrightarrow1
 \label{eq:Q2-final}
\end{equation}
cannot hold.
\end{theorem}

\begin{proof}
Assume \eqref{eq:Q2-final}.  If $F(t)$ denotes the numerator integral, then for every fixed
$0\le z_1<z_2\le1$,
\begin{equation}
 \frac{F(2^k(1+z_2))-F(2^k(1+z_1))}{2^k}\longrightarrow z_2-z_1.
 \label{eq:partial-shell-linearity}
\end{equation}
The exact factorization in the coordinates $x=2^k(1+z)$ gives
\[
 \frac{|\Phi(2^k(1+z))|}{E_{\kappa_1}(2^k(1+z))}
 =\rho_1^{-(k+1)}W(z)\prod_{j=0}^{k}|\sin(\pi2^jz)|,
\]
where $W$ is continuous and positive on $[0,1)$.  Combining this identity with
\eqref{eq:stabilizing-shape} and the Ruelle--Perron--Frobenius limit shows that every
subsequential shell limit in \eqref{eq:partial-shell-linearity} is a constant multiple of
\[
 \nu_1\!\left(\frac{W}{u}\mathbf1_{[z_1,z_2]}\right),
\]
where $\nu_1$ is the eigenmeasure of $\mathcal L_1^*$.  The full-shell case forces the
multiplicative constant to be finite and nonzero.  Equation
\eqref{eq:partial-shell-linearity} would therefore make the measure
$(W/u)\nu_1$ proportional to Lebesgue measure.  Since $W/u$ is continuous and positive on
$[0,1)$ and $\nu_1$ has no endpoint atoms, this would imply $\nu_1\ll\operatorname{Leb}$.

That is impossible.  If $\Differential\nu_1=r\Differential x$ had a nonzero absolutely continuous part, the
eigenmeasure equation would give
\begin{equation}
 |\sin(\pi z)|r(Tz)=\rho_1r(z)
 \quad\text{a.e.}
 \label{eq:conformal-density-final}
\end{equation}
On the positive support of $r$, iteration and Birkhoff's theorem yield
\[
 r(T^nz)=r(z)\frac{\rho_1^n}
 {\prod_{j=0}^{n-1}|\sin(\pi T^jz)|}
 =r(z)\exp\bigl(n(\log(2\rho_1)+o(1))\bigr).
\]
Here $2\rho_1>1$; for example, Cauchy--Schwarz together with the exact $L^2$ identity and
the Gelfond upper bound gives $\rho_1\ge1/\sqrt3>1/2$.  On the other hand, invariance of
Lebesgue measure, Markov's inequality, and Borel--Cantelli imply
$r(T^nz)\le\EulerE^{\varepsilon n}$ eventually for a.e. $z$, for every $\varepsilon>0$.
Choosing $\varepsilon<\log(2\rho_1)$ is a contradiction.  Thus $\nu_1$ is singular and
\eqref{eq:Q2-final} cannot hold.
\end{proof}

The hypothesis covers all gauges of the form
$E_{\kappa_1}(x)L(x)p(\{\log_2x\})$, where $L$ is regularly varying (slow variation is the
most relevant case) and $p$ is positive, continuous, and periodic.  It also covers the
usual log-exp-algebraic or Hardy-field candidates of the correct rough size whenever their
relative logarithmic derivative has a finite limit, because the mean-value theorem then
gives a uniform limiting shell shape.  It does not, by itself, rule out a deliberately
scale-dependent analytic gauge whose within-shell shape keeps changing and attempts to
follow finer and finer dyadic cylinders.  None of the
source materials proves such a fully unrestricted nonexistence theorem.

\subsection{The multiplicatively natural mean}

The endpoint obstruction disappears if additive Cesaro averaging is replaced by logarithmic
averaging.  Each dyadic shell then has equal weight, and the final incomplete shell becomes
negligible.

\begin{theorem}[Unrestricted logarithmic mean]
\label{thm:logarithmic-mean}
There is a constant $A_1^{\log}>0$ such that, for every fixed $t_0>1$,
\begin{equation}
 \frac1{\log(t/t_0)}\int_{t_0}^t
 \frac{|\Phi(x)|}{E_{\kappa_1}(x)}\frac{\Differential x}{x}
 \longrightarrow A_1^{\log}
 \qquad(t\to\infty)
 \label{eq:logarithmic-mean}
\end{equation}
without any restriction on the endpoint.  In the normalization of
\Cref{prop:RPF}, if
$f_1(x)=W_{\kappa_1}((x+1)/2)$, then
\begin{equation}
 A_1^{\log}=\frac1{\log2}
 \left(\int_0^1h_1(x)\Differential x\right)
 \nu_1\!\left(\frac{f_1}{1+\cdot}\right).
 \label{eq:A1log}
\end{equation}
Numerically,
\begin{equation}
 A_1^{\log}=0.09386063575678\ldots .
 \label{eq:A1log-numeric}
\end{equation}
\end{theorem}

\begin{proof}
Let
\[
 \beta_n=\int_{2^{n-1}}^{2^n}
 \frac{|\Phi(x)|}{E_{\kappa_1}(x)}\frac{\Differential x}{x}.
\]
The master factorization and $x=2^ny$, followed by $z=2y-1$, give
\[
 \beta_n=\rho_1^{-n}\int_0^1
 \frac{f_1(z)}{1+z}\prod_{j=0}^{n-1}|\sin(\pi T^jz)|\Differential z.
\]
The spectral limit in \Cref{prop:RPF} yields $\beta_n\to\beta_\infty$, where
$\beta_\infty=(\int h_1)\nu_1(f_1/(1+\cdot))$.  If
$2^N\le t<2^{N+1}$, the logarithmic integral equals
$\sum_{n<N}\beta_n+O(1)$, while $\log(t/t_0)=N\log2+O(1)$.  Ordinary Cesaro convergence of
the sequence $(\beta_n)$ proves \eqref{eq:logarithmic-mean} and
\eqref{eq:A1log}.
\end{proof}

\section{The sharp uniform envelope}
\label{sec:uniform}

The $L^\infty$ numerator rate can be proved by a sharp elementary inequality.  Put
\begin{equation}
 \lambda:=\frac{\sqrt3}{2}.
 \label{eq:lambda}
\end{equation}

\begin{lemma}[Sharp two-step logarithmic inequality]
\label{lem:sharp-sine-inequality}
For every $x\in\RealNumbers$,
\begin{equation}
 \frac23\log s(x)+\frac13\log s(Tx)\le\log\lambda,
 \label{eq:weighted-log-inequality}
\end{equation}
with the usual interpretation at zeros.  Equality holds precisely when
$x\equiv1/3$ or $2/3\pmod1$.
\end{lemma}

\begin{proof}
Cubing the desired inequality gives
\[
 s(x)^2s(Tx)\le\lambda^3.
\]
With $\theta=\pi x$ and $u=\sin^2\theta$,
\[
 s(x)^2s(Tx)=2\AbsoluteValue{\sin^3\theta\cos\theta},
\]
whose square is $4u^3(1-u)$.  The maximum of $u^3(1-u)$ on $[0,1]$ is $27/256$, attained
at $u=3/4$.  Therefore the maximum of the original expression is
$3\sqrt3/8=\lambda^3$, with equality at $x\equiv1/3,2/3$.
\end{proof}

\begin{proposition}[Uniform lacunary-product bound]
\label{prop:Bn-uniform}
For every $n\ge1$,
\begin{equation}
 \lambda^n\le\sup_{x\in[0,1]}
 \prod_{j=0}^{n-1}s(T^jx)\le\lambda^{n-1}.
 \label{eq:Bn-uniform}
\end{equation}
Consequently
\begin{equation}
 \lim_{n\to\infty}\Norm{\prod_{j=0}^{n-1}s\circ T^j}_{\infty}^{1/n}
 =\lambda=\frac{\sqrt3}{2}.
 \label{eq:Gelfond-rate}
\end{equation}
\end{proposition}

\begin{proof}
At $x=1/3$ the orbit alternates between $1/3$ and $2/3$, and every factor equals
$\lambda$, proving the lower bound.  For the upper bound, sum
\eqref{eq:weighted-log-inequality} for $x,Tx,\dots,T^{n-2}x$.  If
$\phi=\log s$, the result is
\[
 3\sum_{j=0}^{n-2}\phi(T^jx)-\phi(x)+\phi(T^{n-1}x)
 \le3(n-1)\log\lambda.
\]
After adding the final factor,
\[
 \sum_{j=0}^{n-1}\phi(T^jx)
 \le(n-1)\log\lambda+\frac13\phi(x)+\frac23\phi(T^{n-1}x)
 \le(n-1)\log\lambda,
\]
because $\phi\le0$.  Exponentiation gives the upper bound.
\end{proof}

This is the classical Gelfond exponent for the Thue--Morse polynomial: by
\eqref{eq:TM-modulus}, the unnormalized polynomial grows like $3^{n/2}$ in supremum norm.
The ergodic-optimization formulation and its generalizations are developed in
\cite{fan-schmeling-shen}.

\begin{theorem}[Sharp global upper envelope]
\label{thm:uniform-envelope}
Define
\begin{equation}
 \boxed{\displaystyle
 \kappa_\infty
 :=\frac{\log\pi+a/2-\log(\sqrt3/2)}{a}
 =\frac32+\frac{\log(\pi/\sqrt3)}{\log2}
 =2.359014879111741\ldots .}
 \label{eq:kappa-infty}
\end{equation}
Then there is a constant $C<\infty$ such that
\begin{equation}
 \AbsoluteValue{\Phiup(x)}\le C E_{\kappa_\infty}(x),
 \qquad x\ge1.
 \label{eq:uniform-envelope}
\end{equation}
This exponent is sharp.  In fact, for every $n\ge0$,
\begin{equation}
 \AbsoluteValue{\Phiup(2^n\cdot2/3)}
 =C_*E_{\kappa_\infty}(2^n\cdot2/3),
 \label{eq:exact-peak-ray}
\end{equation}
where
\begin{equation}
 C_*=W_{\kappa_\infty}(2/3)
 =0.139129774734829\ldots>0.
 \label{eq:Cstar}
\end{equation}
\end{theorem}

\begin{proof}
In \eqref{eq:master-normalization}, the definition of $\kappa_\infty$ makes
\[
 \exp(n(\kappa_\infty a-\log\pi-a/2))=\lambda^{-n}.
\]
After the substitution $x=Ty$, the product is of the form bounded in
\eqref{eq:Bn-uniform}; therefore
\[
 \frac{\AbsoluteValue{\Phiup(2^ny)}}{E_{\kappa_\infty}(2^ny)}
 \le\lambda^{-1}\max_{1/2\le y\le1}W_{\kappa_\infty}(y).
\]
This proves the uniform upper bound.  At $y=2/3$, the orbit alternates between $2/3$ and
$1/3$, so $B_n(2/3)=\lambda^n$ exactly.  Equation
\eqref{eq:master-normalization} then reduces to \eqref{eq:exact-peak-ray}.
\end{proof}

\begin{corollary}[Universal logarithmic order]
\label{cor:universal-log-order}
With the convention $\log0=-\infty$,
\begin{equation}
 \limsup_{x\to\infty}\frac{\log|\Phiup(x)|}{(\log x)^2}
 =-\frac1{2\log2},
 \qquad
 \liminf_{x\to\infty}\frac{\log|\Phiup(x)|}{(\log x)^2}=-\infty.
 \label{eq:universal-log-order}
\end{equation}
Equivalently, the coefficient $1/(2\log2)$ is the sharp global log-square order, but no
positive two-sided pointwise equivalent exists.
\end{corollary}

\begin{proof}
The upper bound in the limsup follows from \eqref{eq:uniform-envelope}; the exact period-two
ray \eqref{eq:exact-peak-ray} gives the reverse inequality.  The liminf is $-\infty$ because
$\Phiup$ vanishes at every positive integer (and also because one may approach those zeros
arbitrarily closely).
\end{proof}

\begin{corollary}[Sharp smoothness class]
\label{cor:superpoly-subexp}
For every $M>0$, $\Phiup(x)=\BigO(x^{-M})$ as $x\to\infty$.  On the other hand, for every
$c>0$, $\EulerE^{cx}\AbsoluteValue{\Phiup(x)}$ is unbounded.  Thus the transform decays faster than every
power but not at any fixed exponential rate.
\end{corollary}

\begin{proof}
The first statement follows from \eqref{eq:uniform-envelope}, because the negative
quadratic in $\log x$ dominates $-M\log x$.  The exact peak ray
\eqref{eq:exact-peak-ray} decays only like $\exp(-\BigO((\log x)^2))$, which is much larger
than $\EulerE^{-cx}$ along that subsequence.
\end{proof}


\section{The raw dyadic-annulus maximum and its Lambert boundary layer}
\label{sec:annular-maxima}

The sharp pointwise envelope of \Cref{thm:uniform-envelope} is evaluated at the actual
argument $x$.  A different question is to compare every value in the shell
$[2^k,2^{k+1}]$ with the envelope at its \emph{left endpoint}.  The maximizer then moves
toward that endpoint and pays a second, slower log-normal cost.  This is the main result of
Wave 1, Document 2 \cite{wave1-doc2}; we reproduce the proof because it is the most delicate
asymptotic in the corpus.

Put
\begin{equation}
 \mathfrak M_k:=\max_{2^k\le x\le2^{k+1}}|\Phiup(x)|,
 \qquad n:=k+1,
 \qquad \beta:=\log\frac{\sqrt3}{2}<0.
 \label{eq:annular-M}
\end{equation}
Also define the tail factor
\[
 \mathcal H(y):=\prod_{m=1}^{\infty}\left|\SincRad\!\left(\frac{\pi y}{2^m}\right)\right|,
 \qquad 1\le y\le2.
\]

\subsection{Exact reduction and transition coordinates}

Writing $x=2^k(1+\theta)$, $0\le\theta\le1$, gives
\begin{equation}
 |\Phiup(2^k(1+\theta))|
 =\frac{2^{-k(k+1)/2}}{\pi^n}
 \mathcal H(1+\theta)(1+\theta)^{-n}Q_n(\theta),
 \label{eq:annular-exact-final}
\end{equation}
where
\[
 Q_n(\theta)=\prod_{j=0}^{n-1}|\sin(\pi2^j\theta)|.
\]
Hence
\begin{equation}
 \mathfrak M_k=\pi^{-n}2^{-k(k+1)/2}A_n,
 \quad
 A_n:=\sup_{0\le\theta\le1}
 \mathcal H(1+\theta)(1+\theta)^{-n}Q_n(\theta).
 \label{eq:An-final}
\end{equation}
The weight $(1+\theta)^{-n}$ pushes the maximum toward $0$, where $Q_n$ vanishes.
For $0<\theta\le1/2$, write uniquely up to endpoints
\begin{equation}
 \theta=2^{-r}s,
 \qquad r\in\IntegerNumbers_{\ge0},\quad s\in I:=[1/4,1/2].
 \label{eq:transition-final}
\end{equation}
If $r<n$, then
\begin{equation}
 Q_n(2^{-r}s)=R_r(s)Q_{n-r}(s),
 \quad
 R_r(s)=\prod_{h=1}^{r}\sin\!\left(\frac{\pi s}{2^h}\right).
 \label{eq:transition-product-final}
\end{equation}
Uniformly for $s\in I$,
\begin{equation}
 \log R_r(s)
 =-\frac a2r^2+r\left(\log(\pi s)-\frac a2\right)+O(1).
 \label{eq:small-block-final}
\end{equation}
The remaining tail satisfies $Q_m(s)\le C\EulerE^{m\beta}$ by
\Cref{prop:Bn-uniform}.

For a matching lower bound one needs near-maximizing tails at essentially arbitrary $s$.

\begin{lemma}[Dense preimages of the maximizing cycle]
\label{lem:dense-preimages-final}
There is $C>0$ such that for every $s_0\in I$ and $q\ge1$ one can choose $s_q\in I$ with
\begin{equation}
 |s_q-s_0|\le2^{-q},
 \qquad T^qs_q\in\{1/3,2/3\},
 \label{eq:preimage-close-final}
\end{equation}
and, for every $m\ge q$,
\begin{equation}
 Q_m(s_q)\ge\exp(m\beta-Cq^2).
 \label{eq:preimage-lower-final}
\end{equation}
\end{lemma}

\begin{proof}
The $q$-fold preimages $(j+1/3)/2^q$ have spacing $2^{-q}$; together with the preimages of
$2/3$ they approximate every point of $I$ as in \eqref{eq:preimage-close-final}.  Before
time $q$, every orbit point is a nonintegral rational with reduced denominator at most
$3\cdot2^{q-j}$.  Therefore
\[
 |\sin(\pi T^js_q)|\ge2\operatorname{dist}(T^js_q,\IntegerNumbers)
 \ge\frac{2}{3\cdot2^{q-j}}.
\]
The first $q$ factors lose at most $\exp(O(q^2))$.  Thereafter the orbit alternates between
$1/3$ and $2/3$, and every factor equals $\sqrt3/2=\EulerE^\beta$.
\end{proof}

Define the finite-dimensional objective
\begin{equation}
 \DecayOptimizationObjective{n}(r,s)
 =-\frac a2r^2+r\left(\log(\pi s)-\frac a2-\beta\right)
 -n\log(1+s2^{-r}).
 \label{eq:Fn-final}
\end{equation}

\begin{lemma}[Reduction to the boundary-layer objective]
\label{lem:An-reduction-final}
As $n\to\infty$,
\begin{equation}
 \log A_n=n\beta+
 \sup_{\substack{r\in\IntegerNumbers_{\ge0}\\s\in I}}\DecayOptimizationObjective{n}(r,s)
 +O((\log\log n)^2).
 \label{eq:An-reduction-final}
\end{equation}
\end{lemma}

\begin{proof}
The upper bound follows from \eqref{eq:small-block-final}, the uniform Gelfond bound, and the
fact that $\theta\ge1/2$ or $r\ge n$ is exponentially inferior to the trial point
$r=\lfloor\log_2n\rfloor$, $s=1/3$.  At every near-maximizer one obtains
$r=O(\log n)$ and $ns2^{-r}=O((\log n)^2)$.  For the lower bound, approximate its $s$ by
$s_q$ from \Cref{lem:dense-preimages-final} with
$q=\lceil4\log_2\log(n+\EulerE^\EulerE)\rceil$.  The derivative of $\DecayOptimizationObjective{n}$ in $s$ is
$O((\log n)^2)$ in the localized region, so the displacement $2^{-q}$ costs $O(1)$; the
tail costs only $O(q^2)=O((\log\log n)^2)$.  The bounded factors
$\mathcal H$ and the finite sinc tail contribute $O(1)$.
\end{proof}

\subsection{Lambert optimization}

Set
\begin{equation}
 C_0:=\log\pi-\frac a2-\beta=\log\!\left(\pi\sqrt{\frac23}\right),
 \qquad
 c:=a\EulerE^{-C_0}=\frac a\pi\sqrt{\frac32}.
 \label{eq:annular-c}
\end{equation}
At the localized maximum, replacing $-n\log(1+s2^{-r})$ by $-ns2^{-r}$ changes the
supremum by $o(1)$.  Put $A=C_0+\log s$ and $z=ar-A$.  Then
\begin{equation}
 -\frac a2r^2+rA-ns2^{-r}
 =-\frac{z^2}{2a}-n\EulerE^{-C_0}\EulerE^{-z}+\frac{A^2}{2a}.
 \label{eq:z-objective}
\end{equation}
As $s$ traverses $[1/4,1/2]$, the interval of possible $A$ has length exactly $a$.
Successive integer values of $r$ therefore make the corresponding $z$-intervals tile the
real line.  Since $A^2/(2a)=O(1)$, the supremum is, up to $O(1)$, the real maximum of
\[
 \phi_n(z)=-\frac{z^2}{2a}-n\EulerE^{-C_0}\EulerE^{-z}.
\]
Its critical equation is $z\EulerE^z=cn$, so the maximizer is $w_n=W(cn)$ and
\begin{equation}
 \max\phi_n=-\frac{w_n^2+2w_n}{2a}.
 \label{eq:Lambert-value}
\end{equation}

\begin{theorem}[Sharp dyadic-annulus maximum]
\label{thm:annular-main}
Let $w_k=W(c(k+1))$, with $c$ from \eqref{eq:annular-c}.  Then
\begin{equation}
 \boxed{\displaystyle
 \log\mathfrak M_k
 =-\frac a2k(k+1)-(k+1)\log\pi+(k+1)\beta
 -\frac{w_k^2+2w_k}{2a}
 +O((\log\log k)^2).}
 \label{eq:annular-main}
\end{equation}
Equivalently,
\begin{equation}
 \log\mathfrak M_k
 =-\frac a2k^2-a\kappa_\infty k
 -\frac{W(c(k+1))^2+2W(c(k+1))}{2a}
 +O((\log\log k)^2).
 \label{eq:annular-main-compact}
\end{equation}
Any maximizing phase $\theta_k$ satisfies
\begin{equation}
 \theta_k=\frac{W(c(k+1))}{a(k+1)}\EulerE^{o(1)}.
 \label{eq:annular-location}
\end{equation}
\end{theorem}

\begin{proof}
Combine \eqref{eq:An-final}, \Cref{lem:An-reduction-final}, and
\eqref{eq:Lambert-value}.  The compact form uses
$a\kappa_\infty=a/2+\log\pi-\beta$.  Strict concavity of $\phi_n$ near $w_n$ and the
$O((\log\log n)^2)$ error localize $z$ to $w_n+o(1)$.  Since
$\theta=s2^{-r}=\EulerE^{-C_0-z}$ and
$n\EulerE^{-C_0-w_n}=w_n/a$, this gives \eqref{eq:annular-location}.
\end{proof}

Let $R=2^k$, $L=\log R$, and $d=\pi^{-1}\sqrt{3/2}$.  Since
$c(k+1)=d(L+a)$, \eqref{eq:annular-main-compact} becomes
\begin{equation}
 \log\max_{R\le x\le2R}|\Phiup(x)|
 =-\frac{L^2}{2a}-\kappa_\infty L
 -\frac{W(d(L+a))^2+2W(d(L+a))}{2a}
 +O((\log\log\log R)^2).
 \label{eq:annular-physical}
\end{equation}
The extra term begins
\begin{equation}
 -\frac{(\log\log R)^2}{2a}
 +\frac{\log\log R\,\log\log\log R}{a}
 +O(\log\log R).
 \label{eq:second-lognormal}
\end{equation}
Thus the annular maximum is smaller than the left-endpoint envelope
$E_{\kappa_\infty}(R)$ by a second log-normal factor, even though the pointwise envelope
itself is attained exactly on the fixed ray $x=2^n(2/3)$.  There is no contradiction: the
fixed ray is a constant fraction into its shell, whereas the raw annular maximizer migrates
toward the left boundary according to \eqref{eq:annular-location}.

\section{Exceptional extrema near dyadic boundaries}
\label{sec:exceptional-extrema}

The sharp envelope in \Cref{thm:uniform-envelope} describes the largest possible values in
a shell, but it does not make every local extremum comparable to one constant.  The deepest
and most explicit obstruction occurs immediately after powers of two.

\begin{theorem}[Dyadic valleys: exact extremal asymptotic]
\label{thm:dyadic-valleys}
Fix an integer $r\ge0$, and put
\[
 b:=r+1,
 \qquad
 d:=1+\TwoAdicValuation(b),
 \qquad
 N:=2^{n-1}.
\]
Let $\xi_{N+r}$ be the unique critical point in $(N+r,N+r+1)$ and
$M_{N+r}=\AbsoluteValue{\Phiup(\xi_{N+r})}$.  Then, as $n\to\infty$,
\begin{align}
 \xi_{N+r}
 &=N+b-\frac{db}{n}+o(n^{-1}),
 \label{eq:valley-location}\\
 M_{N+r}
 &\sim
 \AbsoluteValue{\Phiup(1/2)\Phiup(b/2^d)}
 \left(\frac bN\right)^n
 \left(\frac d{\EulerE n}\right)^d.
 \label{eq:valley-amplitude}
\end{align}
In particular,
\begin{equation}
 \log M_{N+r}
 =-\frac{(\log N)^2}{\log2}+\BigO(\log N).
 \label{eq:valley-log-rate}
\end{equation}
The leading quadratic coefficient in \eqref{eq:valley-log-rate} is twice the coefficient in
$E_\kappa$.
\end{theorem}

\begin{proof}
Write $x=N+z$, $r<z<b$, in the exact boundary identity
\eqref{eq:dyadic-boundary}.  Uniformly for $z\in[r,b]$,
\begin{equation}
 \frac{\AbsoluteValue{\Phiup(\tfrac12+z/(2N))}}
      {\AbsoluteValue{\Phiup(z/(2N))}}
 \longrightarrow\AbsoluteValue{\Phiup(1/2)},
 \qquad
 \left(1+\frac zN\right)^{-n}\longrightarrow1,
 \label{eq:valley-uniform-factors}
\end{equation}
because $n/N\to0$.  Hence
\begin{equation}
 M_{N+r}
 =\AbsoluteValue{\Phiup(1/2)}N^{-n}
   \max_{r\le z\le b}\bigl(z^n\AbsoluteValue{\Phiup(z)}\bigr)(1+o(1)).
 \label{eq:valley-reduced-max}
\end{equation}

By \eqref{eq:zero-local-coefficient}, as $z\uparrow b$,
\[
 \AbsoluteValue{\Phiup(z)}
 =D_b(b-z)^d(1+\BigO(b-z)),
 \qquad
 D_b=b^{-d}\AbsoluteValue{\Phiup(b/2^d)}.
\]
Any fixed region $z\le b-\delta$ is exponentially negligible relative to $b^n$, so the
maximum lies in a shrinking neighborhood of $b$.  Set
$u=n(b-z)/b$.  Uniformly for bounded $u$,
\[
 z^n\AbsoluteValue{\Phiup(z)}
 =D_b b^{n+d}n^{-d}\,u^d\left(1-\frac un\right)^n(1+o(1))
 \longrightarrow
 D_b b^{n+d}n^{-d}u^d\EulerE^{-u}.
\]
The function $u^d\EulerE^{-u}$ has its unique maximum at $u=d$, with value $d^d\EulerE^{-d}$.
Standard localization of the maximum therefore gives
\[
 b-z_{n,r}\sim\frac{db}{n}
\]
and
\[
 \max z^n\AbsoluteValue{\Phiup(z)}
 \sim D_b b^{n+d}n^{-d}d^d\EulerE^{-d}.
\]
Insert $D_b=b^{-d}\AbsoluteValue{\Phiup(b/2^d)}$ into
\eqref{eq:valley-reduced-max} to obtain \eqref{eq:valley-amplitude}; the location statement
is \eqref{eq:valley-location}.  Finally,
$-n\log N=-(\log N)^2/a+\BigO(\log N)$, proving
\eqref{eq:valley-log-rate}.
\end{proof}

For example, in the first interval after $N=2^{n-1}$, $r=0$, $b=d=1$, and
\begin{equation}
 M_N\sim\frac{\Phiup(1/2)^2}{\EulerE n}N^{-n},
 \qquad
 \xi_N=N+1-\frac1n+o(n^{-1}).
 \label{eq:first-dyadic-valley}
\end{equation}
Numerically,
$\Phiup(1/2)=0.553771275888810\ldots$.

\subsection{Consequences for the sequence of extrema}

\begin{corollary}[The upper envelope is sharp but not two-sided]
\label{cor:extrema-limsup-liminf}
For $M_m=\max_{[m,m+1]}\AbsoluteValue{\Phiup}$,
\begin{equation}
 0<\limsup_{m\to\infty}\frac{M_m}{E_{\kappa_\infty}(m)}<\infty,
 \qquad
 \liminf_{m\to\infty}\frac{M_m}{E_{\kappa_\infty}(m)}=0.
 \label{eq:extrema-limsup-liminf}
\end{equation}
Therefore no constant multiple of the sharp uniform envelope makes all local extrema
approach one nonzero amplitude.
\end{corollary}

\begin{proof}
The upper bound follows from \eqref{eq:uniform-envelope}.  Let
$x_n=2^n(2/3)$ and $m_n=\Floor{x_n}$.  The local maximum in the containing unit interval is
at least $\AbsoluteValue{\Phiup(x_n)}$, and \eqref{eq:exact-peak-ray} gives a positive lower bound
after division by $E_{\kappa_\infty}(m_n)$; changing the argument by less than one changes
$E_{\kappa_\infty}$ by a factor tending to one.  This proves the positive limsup.

For the liminf, take $m=N+r$ with fixed $r$.  Comparing
\eqref{eq:valley-log-rate} with
$\log E_{\kappa_\infty}(N)=-(\log N)^2/(2a)+\BigO(\log N)$ shows that the quotient tends to
zero.
\end{proof}

This is a stronger obstruction than the integer zeros alone: even after replacing the
function by the absolute values of its unique local extrema, a constant-amplitude
normalization does not exist.

\section{A continuum of intermediate boundary rates}
\label{sec:intermediate-rates}

The boundary identity also explains why many intermediate effective decay rates appear.
Suppose $N=2^{n-1}$, $z=z_n=o(N)$, and $z_n\to\infty$.  From
\eqref{eq:dyadic-boundary}, provided the tail ratio remains bounded away from zero and
infinity,
\begin{equation}
 \log\AbsoluteValue{\Phiup(N+z_n)}
 =\log\AbsoluteValue{\Phiup(z_n)}+n\log\frac{z_n}{N}+o(n).
 \label{eq:intermediate-master}
\end{equation}
If $\log z_n\sim\alpha\log N$ with $0\le\alpha<1$ and $z_n$ itself follows a typical
fixed-phase decay law, then substitution of \eqref{eq:typical-ray} into
\eqref{eq:intermediate-master} gives the quadratic coefficient
\begin{equation}
 \log\AbsoluteValue{\Phiup(N+z_n)}
 =-\frac{1-\alpha+\alpha^2/2}{a}(\log N)^2
  +\BigO(\log N)+o((\log N)^2).
 \label{eq:intermediate-quadratic}
\end{equation}
The coefficient decreases continuously from $1/a$ at $\alpha=0$ to $1/(2a)$ as
$\alpha\uparrow1$.  This calculation is conditional only on the specified decay law for
$z_n$ and makes precise the recursive origin of the broad scatter visible in plots of the
extrema.

\section{Numerical evaluation of the finite-\texorpdfstring{$q$}{q} spectrum}
\label{sec:numerics}

\subsection{Chebyshev collocation}

For smooth weights (in particular $q\ge1$), the Perron root of
\eqref{eq:transfer-operator} is efficiently obtained by Chebyshev--Lobatto collocation.
Let
\[
 x_j=\frac{1-\cos(j\pi/(N-1))}{2},\qquad 0\le j<N,
\]
interpolate $f$ at the points $x_j/2$ and $(x_j+1)/2$ by the barycentric formula, and form
the matrix representation of $\mathcal L_q$.  Its leading eigenvalue converges rapidly to
$\rho_q$.

For $q=1$ the convergence is:
\begin{center}
\begin{tabular}{@{}rr@{}}
\toprule
$N$ & leading collocation eigenvalue\\
\midrule
8  & 0.661322639955074\\
10 & 0.661322602233874\\
12 & 0.661322602061258\\
16 & 0.661322602060566\\
24 & 0.661322602060565\\
\bottomrule
\end{tabular}
\end{center}
The rational enclosure in \Cref{prop:rho1-rigorous-bound} is independent of this numerical
convergence and provides a rigorous check on the computed value.

\subsection{Sample spectrum}

The following values illustrate the interpolation from typical behavior to the sharp
uniform envelope.  The rows $q=0,2,\infty$ are exact; the remaining rows are numerical
collocation values.

\begin{center}
\begin{tabular}{@{}rlll@{}}
\toprule
$q$ & $\rho_q$ & $\Lambda_q=\rho_q^{1/q}$ & $\kappa_q$\\
\midrule
$0$ & -- & $0.5000000000$ & $3.1514961295$\\
$1$ & $0.6613226021$ & $0.6613226021$ & $2.7480700149$\\
$2$ & $0.5000000000$ & $0.7071067812$ & $2.6514961295$\\
$3$ & $0.3948966324$ & $0.7336593838$ & $2.5983138062$\\
$4$ & $0.3201941016$ & $0.7522346457$ & $2.5622414703$\\
$8$ & $0.1612444589$ & $0.7960412993$ & $2.4805809434$\\
$16$& $0.0500719357$ & $0.8293247927$ & $2.4214870020$\\
$\infty$ & -- & $0.8660254038$ & $2.3590148791$\\
\bottomrule
\end{tabular}
\end{center}

\begin{center}
\begin{tikzpicture}
\begin{axis}[
 width=0.85\textwidth,
 height=0.42\textwidth,
 xlabel={$q/(1+q)$},
 ylabel={$\kappa_q$},
 xmin=0,xmax=1,
 ymin=2.30,ymax=3.22,
 grid=major,
 tick align=outside,
]
\addplot+[mark=*] coordinates {
 (0,3.1514961295)
 (0.5,2.7480700149)
 (0.6666666667,2.6514961295)
 (0.75,2.5983138062)
 (0.8,2.5622414703)
 (0.8888888889,2.4805809434)
 (0.9411764706,2.4214870020)
 (1,2.3590148791)
};
\end{axis}
\end{tikzpicture}
\end{center}

The monotone curve is not a cosmetic refinement: choosing one row answers one specific
notion of ``amplitude.''  Moving from $q=0$ to $q=\infty$ changes the power of $x$ by almost
$0.8$, an effect much larger than any logarithmic correction.

\section{Source-by-source reconciliation and correction of the draft}
\label{sec:source-reconciliation}

The preceding theorems now permit a precise adjudication of every disputed claim.

\subsection{What all four wave-1 documents got right}

All four identify the universal factor
\[
 \exp\!\left(-\frac{(\log x)^2}{2\log2}\right)
\]
and reduce the residual behavior to a dyadic lacunary product.  They also agree on the
integer zero divisor, the Thue--Morse sign law, the existence of one lobe maximum per unit
interval, the typical exponent $\kappa_0$, and the sharp pointwise-envelope exponent
$\kappa_\infty$.  These are all correct.

\subsection{Why the proposed answers appeared to disagree}

Documents 1 and 2 primarily study extremal quantities; Documents 3 and 4 primarily study
means.  The power of $x$ is determined by the numerator scale being sampled:
\[
 \frac12\quad(q=0),\qquad
 \rho_1\quad(q=1),\qquad
 \frac1{\sqrt2}\quad(q=2),\qquad
 \frac{\sqrt3}{2}\quad(q=\infty).
\]
Substituting these four scales into
$\kappa=\tfrac12+\log_2(\pi/\Lambda)$ gives four different, simultaneously correct
answers.  The old draft's value $\log_2(2\pi)$ is exactly the $q=2$ row.

Document 2's Lambert term is not a fifth competing power.  It refines the $q=\infty$ row
when the envelope is anchored at the left endpoint of an entire dyadic annulus.  Likewise,
Document 4's fixed-offset valleys do not contradict the upper envelope: they show that the
sequence of all local maxima has zero liminf on that upper scale.

\subsection{The genuine error surface}

The only substantive numerical contradiction in the corpus is the LIL constant.  Document 1
has $\pi/\sqrt2$ with denominator $\sqrt{n\log\log n}$; Document 3 quotes $\pi$ from
\cite{ahl2017}.  The exact finite-$n$ variance \eqref{eq:variance-exact} and the explicit
martingale decomposition \eqref{eq:Gordin-exact} decide the issue in favor of Document 1.
The published equation misses the factor $1/2$ in the norm of each real cosine mode.

\subsection{The literal Cesaro question}

The weak condition from the original question is solved by
\[
 g(x)=A_1x^{-\kappa_1}
 \exp\!\left(-\frac{(\log x)^2}{2\log2}\right),
 \qquad A_1=0.091266124131588\ldots .
\]
It is positive and analytic on $(0,\infty)$, eventually decreasing, and for every fixed
$m$, $(-1)^mg^{(m)}(x)>0$ for all sufficiently large $x$ (the threshold may depend on $m$).
The normalized ordinary Cesaro average is bounded between positive constants and converges
to one along dyadic endpoints.  For unrestricted endpoints it converges to a nonconstant
mantissa profile.  The stabilizing-gauge theorem proves that no ordinary power of $x$ or
$\log x$, no slowly varying correction, and no fixed continuous log-periodic modulation can
flatten that profile.  The logarithmic average does converge for all endpoints.

The strongest possible wording must remain disciplined: the available proof excludes the
stabilizing class, not every imaginable analytic monotone function whose relative shape
changes at every scale.  The comparative audit's broad phrase ``unattainable in principle''
should therefore be read with this stated hypothesis.

\subsection{No universal logarithmic power}

At every fixed finite $q$, Ruelle--Perron--Frobenius asymptotics give a pure exponential in
the dyadic level and a nonzero constant; multiplying by $(\log x)^p$ forces the normalized
mean to $0$ or $\infty$ unless $p=0$.  At $q=2$ this is already exact by Parseval.  At
$q=\infty$ the period-two ray gives an exact constant ratio.  Powers of the dyadic level do
occur in the exceptional valley formula, but their exponent is the integer zero
multiplicity $1+\nu_2(r+1)$ and they accompany a different quadratic decay coefficient.
They are not a universal half-logarithmic correction.

\subsection{Imported analytic inputs and certified parts}

The elementary identities, variance computation, Gelfond inequality, Lambert optimization,
and fixed-offset valley asymptotic are proved here from scratch.  The exact $L^2$ row is
finite Fourier algebra.  The $L^1$ exponential asymptotic and the enclosure
$0.66130<\rho_1<0.66135$ are published inputs from Fouvry--Mauduit and
Aistleitner--Hofer--Larcher \cite{fouvrymauduit1996,ahl2017}.  The shell-profile limits use
the standard Ruelle--Perron--Frobenius spectral gap for the explicit operator
\eqref{eq:transfer-operator}; this is the one substantial analytic block imported rather
than reproved.  The independent exact Collatz--Wielandt/Sturm certificate in
\Cref{app:certificate} supplies a machine-checkable enclosure for $\rho_1$ without relying
on floating-point collocation.

\section{Conclusion}
\label{sec:conclusion}

The Fourier image of Rvachev's up-function has one deterministic core and several genuinely
different stochastic or extremal refinements.  The exact identity
\[
 |\Phiup(2^ny)|
 =|\Phiup(y)|\pi^{-n}y^{-n}2^{-n(n+1)/2}
 \prod_{j=1}^n|\sin(\pi2^jy)|
\]
separates them perfectly.  The triangular denominator produces
$-(\log x)^2/(2\log2)$.  The lacunary numerator then selects a pressure level:
Lebesgue geometric mean, Perron $L^1$ mean, Parseval $L^2$ mean, or maximizing orbit.

The final quantitative dictionary is:
\begin{align*}
 \kappa_0&=3.151496129472319\ldots &&\text{typical / geometric mean},\\
 \kappa_1&=2.748070014871334\ldots &&\text{arithmetic mean and Cesaro gauge},\\
 \kappa_2&=2.651496129472319\ldots &&\text{RMS; the old draft's exponent},\\
 \kappa_\infty&=2.359014879111741\ldots &&\text{sharp pointwise upper envelope}.
\end{align*}
Typical logarithms fluctuate on the
$\sqrt{\log x\,\log\log\log x}$ scale with corrected constant $\pi/\sqrt2$.
Raw dyadic-annulus maxima carry an additional Lambert-$W$ term of order
$-(\log\log x)^2/(2\log2)$.  Fixed-offset lobe maxima after powers of two can instead have
the doubled leading coefficient $-(\log x)^2/\log2$.

For the question that originally motivated the investigation, the positive analytic gauge
$A_1E_{\kappa_1}$ gives the requested two-sided bounded Cesaro normalization and the exact
dyadic-endpoint limit.  The failure of an unrestricted ordinary limit for all standard
stabilizing gauges is structural: the last dyadic shell retains a singular distribution of
mass.  Logarithmic averaging respects the multiplicative self-similarity and restores an
unrestricted limit.

No substantive disagreement among the source documents remains after the norm, endpoint
convention, and LIL normalization are made explicit.  What looked like one irregular
asymptotic is in fact a compact thermodynamic spectrum, supplemented by one Lambert boundary
layer and one arithmetic valley mechanism.

\appendix

\section{Exact Collatz--Wielandt certificate for \texorpdfstring{$\rho_1$}{rho1}}
\label{app:certificate}

The following SymPy script uses only exact rational arithmetic.  It constructs the rational
function in the proof of \Cref{prop:rho1-rigorous-bound}; Sturm's theorem, invoked by
\texttt{count\_roots}, certifies that the relevant numerators and denominators have no roots
on $[0,1]$.  Positivity at $t=0$ therefore proves positivity throughout the interval.

\begin{lstlisting}[style=code,language=Python]
import sympy as sp

t = sp.symbols("t", real=True)
u = 2*t/(1+t**2)
v = (1-t**2)/(1+t**2)

c = sp.Rational(376189, 10**6)
d = sp.Rational(-69093, 10**6)
e = sp.Rational(15483, 10**6)
F = lambda z: 1 + c*z + d*z**2 + e*z**3

R = sp.cancel((u*F(u) + v*F(v)) / (2*F(2*u*v)))
lo = sp.Rational(66126807, 10**8)
hi = sp.Rational(66134891, 10**8)

for expr in (R - lo, hi - R):
    num, den = map(lambda p: sp.Poly(p, t),
                   sp.together(expr).as_numer_denom())
    assert num.eval(0) > 0 and den.eval(0) > 0
    assert num.count_roots(0, 1) == 0
    assert den.count_roots(0, 1) == 0

print("Certified:", lo, "< rho_1 <", hi)
\end{lstlisting}

For the high-precision decimal, one may replace the exact certificate by a Chebyshev
collocation matrix.  The rapid stabilization displayed in \Cref{sec:numerics} is typical
for this analytic weighted-composition operator.

\section{Notation cross-reference}
\label{app:notation}

\begin{center}
\begin{tabularx}{\textwidth}{@{}lX@{}}
\toprule
Symbol & Meaning\\
\midrule
$\Phiup=\widehat{\RvachevUp}$ & Fourier transform of Rvachev's up-function\\
$a$ & $\log2$\\
$T$ & doubling map $x\mapsto2x\pmod1$\\
$s(x)$ & $\AbsoluteValue{\sin(\pi x)}$\\
$B_n(y)$ & $\prod_{j=1}^n s(T^jy)$\\
$E_\kappa(x)$ & $x^{-\kappa}\exp(- (\log x)^2/(2\log2))$\\
$W_\kappa(y)$ & bounded dyadic phase factor in \eqref{eq:master-normalization}\\
$\mathcal L_q$ & weighted Perron operator \eqref{eq:transfer-operator}\\
$\rho_q$ & Perron root of $\mathcal L_q$\\
$\Lambda_q$ & $\rho_q^{1/q}$, the $L^q$ numerator growth per dyadic level\\
$\kappa_q$ & sharp power in the normalized $L^q$ shell decay\\
$M_m$ & unique maximum of $\AbsoluteValue{\Phiup}$ on $[m,m+1]$\\
\bottomrule
\end{tabularx}
\end{center}

\begin{thebibliography}{99}

\bibitem{arias1982}
J. Arias de Reyna,
\emph{An infinitely differentiable function with compact support: Definition and properties},
English translation of \emph{Definici\'on y estudio de una funci\'on indefinidamente
diferenciable de soporte compacto}, Rev. Real Acad. Cienc. Exact. F\'is. Natur. Madrid
\textbf{76} (1982), 21--38; arXiv:1702.05442 (2017).

\bibitem{arias2018}
J. Arias de Reyna,
\emph{Arithmetic of the Fabius function},
Integers \textbf{18} (2018), Paper A51; arXiv:1702.06487.

\bibitem{baladi}
V. Baladi,
\emph{Positive Transfer Operators and Decay of Correlations},
Advanced Series in Nonlinear Dynamics, vol.~16, World Scientific, 2000.

\bibitem{fabius}
J. Fabius,
\emph{A probabilistic example of a nowhere analytic $C^\infty$-function},
Z. Wahrscheinlichkeitstheorie verw. Gebiete \textbf{5} (1966), 173--174,
\href{https://doi.org/10.1007/BF00536652}{doi:10.1007/BF00536652}.

\bibitem{fan1997}
A.-H. Fan,
\emph{Multifractal analysis of infinite products},
J. Statist. Phys. \textbf{86} (1997), 1313--1336,
\href{https://doi.org/10.1007/BF02183625}{doi:10.1007/BF02183625}.

\bibitem{fan-schmeling-shen}
A. Fan, J. Schmeling, and W. Shen,
\emph{$L^\infty$-estimation of generalized Thue--Morse trigonometric polynomials and
ergodic maximization},
Discrete Contin. Dyn. Syst. \textbf{41} (2021), 297--327,
\href{https://doi.org/10.3934/dcds.2020363}{doi:10.3934/dcds.2020363}.

\bibitem{gelfond}
A. O. Gel'fond,
\emph{Sur les nombres qui ont des propri\'et\'es additives et multiplicatives donn\'ees},
Acta Arith. \textbf{13} (1967/68), 259--265,
\href{https://doi.org/10.4064/aa-13-3-259-265}{doi:10.4064/aa-13-3-259-265}.

\bibitem{jessen-wintner}
B. Jessen and A. Wintner,
\emph{Distribution functions and the Riemann zeta function},
Trans. Amer. Math. Soc. \textbf{38} (1935), 48--88.

\bibitem{mse-extrema}
V. Reshetnikov,
\emph{Extrema of an infinite product of sinc functions},
Mathematics Stack Exchange, question 2742261 (2018),
\url{https://math.stackexchange.com/questions/2742261/}.

\bibitem{mse-decay}
V. Reshetnikov,
\emph{Asymptotic decay rate of an infinite product of sinc functions},
Mathematics Stack Exchange, question 2745497 (2018),
\url{https://math.stackexchange.com/questions/2745497/}.

\bibitem{proveit}
V. Reshetnikov,
\emph{ProveIt: Analysis/FabiusFunction}, GitHub repository,
\url{https://github.com/VladimirReshetnikov/ProveIt/tree/main/Analysis/FabiusFunction}.

\bibitem{decay-draft}
V. Reshetnikov,
\emph{Asymptotic decay rate of an infinite product of sinc functions}, source draft in the
\texttt{ProveIt} repository,
\url{https://github.com/VladimirReshetnikov/ProveIt/blob/main/Analysis/FabiusFunction/docs/semi-formalized-research-frontiers/drafts/rvachev_up_fourier_decay/asymptotic-decay-rate-of-an-infinite-product-of-sinc-functions/asymptotic-decay-rate-of-an-infinite-product-of-sinc-functions.tex}.

\bibitem{queffelec}
M. Queff\'elec,
\emph{Questions around the Thue--Morse sequence},
Unif. Distrib. Theory \textbf{13} (2018), no.~1, 1--25,
\href{https://doi.org/10.1515/udt-2018-0001}{doi:10.1515/udt-2018-0001}.

\bibitem{rvachev1971}
V. L. Rvachev and V. A. Rvachev,
\emph{On one compactly supported function},
Dokl. Akad. Nauk Ukrainian SSR, Ser.~A, no.~8 (1971), 705--707
(in Ukrainian).


\bibitem{ahl2017}
C. Aistleitner, R. Hofer, and G. Larcher,
\emph{On evil Kronecker sequences and lacunary trigonometric products},
Ann. Inst. Fourier (Grenoble) \textbf{67} (2017), no.~2, 637--687,
\href{https://doi.org/10.5802/aif.3094}{doi:10.5802/aif.3094};
preprint arXiv:1502.06738.

\bibitem{fouvrymauduit1996}
\'E. Fouvry and C. Mauduit,
\emph{Sommes des chiffres et nombres presque premiers},
Math. Ann. \textbf{305} (1996), no.~3, 571--600,
\href{https://doi.org/10.1007/BF01444238}{doi:10.1007/BF01444238}.

\bibitem{stout1970}
W. F. Stout,
\emph{A martingale analogue of Kolmogorov's law of the iterated logarithm},
Z. Wahrscheinlichkeitstheorie verw. Gebiete \textbf{15} (1970), 279--290,
\href{https://doi.org/10.1007/BF00533299}{doi:10.1007/BF00533299}.

\bibitem{heyde-scott1973}
C. C. Heyde and D. J. Scott,
\emph{Invariance principles for the law of the iterated logarithm for martingales and
processes with stationary increments},
Ann. Probab. \textbf{1} (1973), 428--436.

\bibitem{wave1-doc1}
\emph{Rvachev up Fourier decay, Wave 1, Document 1}, working report in the
\texttt{ProveIt} repository,
\url{https://github.com/VladimirReshetnikov/ProveIt/blob/main/Analysis/FabiusFunction/docs/semi-formalized-research-frontiers/drafts/rvachev_up_fourier_decay/rvachev_up_fourier_decay-1/rvachev_up_fourier_decay.tex}.

\bibitem{wave1-doc2}
\emph{Rvachev up Fourier decay, Wave 1, Document 2}, working report in the
\texttt{ProveIt} repository,
\url{https://github.com/VladimirReshetnikov/ProveIt/blob/main/Analysis/FabiusFunction/docs/semi-formalized-research-frontiers/drafts/rvachev_up_fourier_decay/Rvachev_Up_Fourier_Decay-2/Rvachev_Up_Fourier_Decay.tex}.

\bibitem{wave1-doc3}
\emph{Fourier decay of Rvachev up, Wave 1, Document 3}, working report in the
\texttt{ProveIt} repository,
\url{https://github.com/VladimirReshetnikov/ProveIt/blob/main/Analysis/FabiusFunction/docs/semi-formalized-research-frontiers/drafts/rvachev_up_fourier_decay/Rvachev_Up_Fourier_Decay-3/Fourier_decay_of_Rvachev_up.tex}.

\bibitem{wave1-doc4}
\emph{Rvachev up Fourier decay, Wave 1, Document 4}, working report in the
\texttt{ProveIt} repository,
\url{https://github.com/VladimirReshetnikov/ProveIt/blob/main/Analysis/FabiusFunction/docs/semi-formalized-research-frontiers/drafts/rvachev_up_fourier_decay/rvachev_up_fourier_decay-4/rvachev_up_fourier_decay.tex}.

\bibitem{comparative-audit}
\emph{Rvachev up Fourier decay comparative audit}, working report in the
\texttt{ProveIt} repository,
\url{https://github.com/VladimirReshetnikov/ProveIt/blob/main/Analysis/FabiusFunction/docs/semi-formalized-research-frontiers/drafts/rvachev_up_fourier_decay/Rvachev_Up_Fourier_Decay_Comparative_Audit/Rvachev_Up_Fourier_Decay_Comparative_Audit.tex}.

\end{thebibliography}

\end{document}
